\documentclass[12pt,a4paper,UTF8]{ctexart}
\usepackage[a4paper,top=22mm,bottom=22mm,left=22mm,right=22mm,headheight=15pt]{geometry}
\newcommand{\Needspace}[1]{\par\begingroup\dimen0=#1\relax\dimen2=\pagegoal\advance\dimen2 by-\pagetotal\ifdim\dimen0>\dimen2\newpage\fi\endgroup}
\newcommand{\needspace}[1]{\par\begingroup\skip0=\lastskip\vskip-\skip0\dimen0=#1\relax\advance\dimen0 by\skip0\vskip\dimen0\penalty-100\vskip-\dimen0\vskip\skip0\penalty9999\endgroup}
\usepackage{amsmath,amssymb,mathtools,bm}
\usepackage{graphicx,xcolor,booktabs,longtable,array,enumitem,fancyhdr,caption,float,microtype,setspace}
\usepackage[hidelinks,unicode]{hyperref}
\usepackage[most]{tcolorbox}
\definecolor{ContestBlue}{HTML}{17365D}
\hypersetup{pdftitle={第066套},pdfauthor={}}
\setstretch{1.24}
\setlength{\parindent}{2em}
\setlength{\parskip}{0.35em}
\setlength{\emergencystretch}{2em}
\setlist{itemsep=0.35em,topsep=0.4em}
\captionsetup{font=small,labelfont=bf,justification=centering}
\setcounter{secnumdepth}{0}
\pagestyle{fancy}
\fancyhf{}
\fancyhead[L]{\small 第066套}
\fancyfoot[C]{\thepage}
\renewcommand{\headrulewidth}{0.4pt}
\ctexset{
  section={format=\Large\bfseries\color{ContestBlue},beforeskip=1.2em,afterskip=0.7em},
  subsection={format=\large\bfseries\color{ContestBlue},beforeskip=1.0em,afterskip=0.5em},
  subsubsection={format=\normalsize\bfseries,beforeskip=0.8em,afterskip=0.35em}
}
\usepackage{calc}
\newcounter{none}
\let\originalsubsection\subsection
\renewcommand{\subsection}{\needspace{6\baselineskip}\originalsubsection}
\makeatletter
\newcommand{\tightlist}{\setlength{\itemsep}{0pt}\setlength{\parskip}{0pt}}
\newcommand{\maxwidth}{\ifdim\Gin@nat@width>\linewidth\linewidth\else\Gin@nat@width\fi}
\newcommand{\maxheight}{\ifdim\Gin@nat@height>0.72\textheight 0.72\textheight\else\Gin@nat@height\fi}
\newcommand{\pandocbounded}[1]{\begingroup\setkeys{Gin}{width=\maxwidth,height=\maxheight,keepaspectratio}#1\endgroup}
\makeatother
\begin{document}
\begin{titlepage}
  \centering
  \vspace*{58mm}
  {\Huge\bfseries\color{ContestBlue} 第066套\par}
  \vfill
  \vspace*{28mm}
\end{titlepage}
\section{第 1 题}\label{ux7b2c-1-ux9898}

一种弹性粗绳被拉伸时密度保持不变，因此拉伸后横截面积会减小，但材料的杨氏模量不变。弹性绳未拉伸时的横截面积为
\(A_0\)、长度为 \(l_0\)；用拉力 \(F\)
将它均匀拉伸后，横截面积和长度分别变为 \(A\) 和 \(l\)，并满足

\[
\frac{F}{A}=Y\frac{l-l_0}{l_0},
\]

其中 \(Y\) 为材料的杨氏模量。

（1）当弹性绳被均匀拉伸到长度 \(l_1\) 时，求绳内拉力 \(F_1\) 和弹性势能
\(E_1\)。

（2）弹性绳的质量为 \(M\)。将绳的一端固定，另一端悬挂质量为 \(m\)
的物体，重力加速度为 \(g\)，求静止时的绳长 \(l_2\)。答案用
\(M,m,g,l_0,A_0,Y\) 表示。

\begin{figure}[H]
\centering
\includegraphics[width=0.685\linewidth,height=\textheight,keepaspectratio,alt={图1-1　弹性绳均匀拉伸与竖直悬挂时的几何量}]{../figures/question-01.pdf}
\caption*{图1-1　弹性绳均匀拉伸与竖直悬挂时的几何量}
\end{figure}

\subsection{一、题型定位}\label{ux4e00ux9898ux578bux5b9aux4f4d}

绳子伸长后会变细。先把体积守恒代入杨氏模量关系，再沿未伸长的材料坐标积分悬绳；不能把形变后的线密度当成常数。

\subsection{二、知识点}\label{ux4e8cux77e5ux8bc6ux70b9}

\begin{enumerate}
\def\labelenumi{\arabic{enumi}.}
\tightlist
\item
  密度不变且质量不变，意味着总体积不变，即 \(Al=A_0l_0\)。
\item
  题给应力---应变关系中的应力使用形变后的横截面积 \(A\)。
\item
  变力做功必须积分，不能把末态拉力直接乘伸长量。
\item
  悬绳各截面承受的下方重量不同，要沿未形变坐标逐段积分。
\item
  有限伸长比为正要求张力小于极限值 \(A_0Y\)。
\end{enumerate}

\subsection{三、物理图像}\label{ux4e09ux7269ux7406ux56feux50cf}

随着绳长增加，体积守恒使横截面积按 \(A=A_0l_0/l\)
减小。因此即使材料仍服从线性应力---应变关系，整根绳的拉力也不再与伸长量严格成正比，而是逐渐趋近
\(A_0Y\)。悬挂时，上端截面承担``物体加整根绳''的重量，伸长最大；下端截面只承担物体重量，伸长较小。

\subsection{四、解题入口}\label{ux56dbux89e3ux9898ux5165ux53e3}

先把体积约束代入本构关系，得到只含当前长度的
\(F(l)\)。第二问选取绳在未伸长状态中的材料坐标
\(\lambda\)，从下端向上量；位于 \(\lambda\) 处的截面承受质量
\(m+M\lambda/l_0\)，再把该微元的伸长比积分。

\subsection{五、解析}\label{ux4e94ux89e3ux6790}

\begin{figure}[H]
\centering
\includegraphics[width=0.705\linewidth,height=\textheight,keepaspectratio,alt={图1-2　以未形变坐标标记悬绳微元及其下方负载}]{../figures/solution-01.pdf}
\caption*{图1-2　以未形变坐标标记悬绳微元及其下方负载}
\end{figure}

图1-2区分材料坐标 \(\lambda\) 与伸长后的空间距离
\(s(\lambda)\)。悬点以下所含材料质量决定拉力，真实截面再由不可压缩条件求出。

\needspace{3\baselineskip}
\textbf{（一）由体积守恒消去横截面积}

绳的密度不变、质量不变，所以拉伸前后体积相同：

\[
Al=A_0l_0. \tag{1}
\]

故

\[
A=\frac{A_0l_0}{l}. \tag{2}
\]

将式（2）代入题给本构关系：

\[
\frac{F}{A_0l_0/l}
=Y\frac{l-l_0}{l_0}. \tag{3}
\]

逐步整理：

\[
\begin{aligned}
F
&=\frac{A_0l_0}{l}\,Y\frac{l-l_0}{l_0}\\
&=A_0Y\frac{l-l_0}{l}\\
&=A_0Y\left(1-\frac{l_0}{l}\right).
\end{aligned} \tag{4}
\]

因此绳长为 \(l_1\) 时

\[
\boxed{F_1=A_0Y\left(1-\frac{l_0}{l_1}\right)}. \tag{5}
\]

当 \(l_1=l_0\) 时 \(F_1=0\)；当 \(l_1\to\infty\) 时
\(F_1\to A_0Y\)，这与式（4）的有限上限一致。

\needspace{3\baselineskip}
\textbf{（二）变力做功求弹性势能}

准静态拉伸中，外力所做的功全部转化为弹性势能。取未伸长状态势能为零：

\[
E_1=\int_{l_0}^{l_1}F(l)\,dl. \tag{6}
\]

代入式（4）：

\[
\begin{aligned}
E_1
&=A_0Y\int_{l_0}^{l_1}\left(1-\frac{l_0}{l}\right)dl\\
&=A_0Y\left[
l-l_0\ln l
\right]_{l_0}^{l_1}\\
&=A_0Y\left[
(l_1-l_0)-l_0\bigl(\ln l_1-\ln l_0\bigr)
\right]\\
&=A_0Y\left[
l_1-l_0-l_0\ln\frac{l_1}{l_0}
\right].
\end{aligned} \tag{7}
\]

所以

\[
\boxed{
E_1=A_0Y\left(l_1-l_0-l_0\ln\frac{l_1}{l_0}\right)
}. \tag{8}
\]

令 \(l_1=l_0+\Delta l\) 且 \(\Delta l\ll l_0\)。利用

\[
\ln(1+x)=x-\frac{x^2}{2}+O(x^3),
\]

可得

\[
E_1\simeq \frac{A_0Y}{2l_0}(\Delta l)^2,
\]

即回到小形变弹簧的势能形式，等效劲度系数为 \(A_0Y/l_0\)。

\needspace{3\baselineskip}
\textbf{（三）建立悬绳微元的张力}

用 \(\lambda\)
表示微元在未形变绳上的坐标，从下端向上量，\(0\le\lambda\le l_0\)。均匀绳的线密度（以未形变长度计）为

\[
\lambda_m=\frac{M}{l_0}. \tag{9}
\]

位于 \(\lambda\) 处的截面下方包括物体 \(m\) 和未形变长度为 \(\lambda\)
的绳段。因此静力平衡给出

\[
T(\lambda)
=\left(m+\frac{M}{l_0}\lambda\right)g. \tag{10}
\]

特别地，

\[
T(0)=mg,\qquad T(l_0)=(m+M)g. \tag{11}
\]

\needspace{3\baselineskip}
\textbf{（四）求单个材料微元的伸长比}

未形变长度为 \(d\lambda\) 的微元，形变后长度记为 \(ds\)，伸长比为

\[
q=\frac{ds}{d\lambda}. \tag{12}
\]

微元体积守恒：

\[
A\,ds=A_0\,d\lambda,
\]

所以

\[
A=\frac{A_0}{q}. \tag{13}
\]

本构关系写成微元形式：

\[
\frac{T}{A}=Y(q-1). \tag{14}
\]

将式（13）代入：

\[
\frac{Tq}{A_0}=Y(q-1). \tag{15}
\]

移项并解出 \(q\)：

\[
\begin{aligned}
A_0Yq-Tq&=A_0Y,\\
q(A_0Y-T)&=A_0Y,\\
q&=\frac{A_0Y}{A_0Y-T}.
\end{aligned} \tag{16}
\]

因此

\[
ds=\frac{A_0Y}{A_0Y-T(\lambda)}\,d\lambda. \tag{17}
\]

\needspace{3\baselineskip}
\textbf{（五）沿整根绳积分}

把式（10）代入式（17）：

\[
l_2
=\int_0^{l_0}
\frac{A_0Y}
{A_0Y-mg-\dfrac{Mg}{l_0}\lambda}
\,d\lambda. \tag{18}
\]

令

\[
u=A_0Y-mg-\frac{Mg}{l_0}\lambda,
\qquad
du=-\frac{Mg}{l_0}d\lambda. \tag{19}
\]

积分上下限为

\[
\lambda=0\Rightarrow u=A_0Y-mg,
\]

\[
\lambda=l_0\Rightarrow u=A_0Y-mg-Mg.
\]

故

\[
\begin{aligned}
l_2
&=-\frac{A_0Yl_0}{Mg}
\int_{A_0Y-mg}^{A_0Y-mg-Mg}\frac{du}{u}\\
&=\frac{A_0Yl_0}{Mg}
\ln\frac{A_0Y-mg}{A_0Y-mg-Mg}.
\end{aligned} \tag{20}
\]

最终

\[
\boxed{
l_2=
\frac{A_0Yl_0}{Mg}
\ln\frac{A_0Y-mg}{A_0Y-(m+M)g}
}. \tag{21}
\]

解存在的条件是

\[
\boxed{A_0Y>(m+M)g}. \tag{22}
\]

这是因为上端张力最大，若它达到或超过
\(A_0Y\)，式（16）的伸长比便不再是有限正数。

\needspace{3\baselineskip}
\textbf{补充推导与适用条件}

令整根绳的伸长比为

\[
q=\frac{l}{l_0}>1.
\]

体积守恒给 \(A=A_0/q\)，题设本构关系于是化为

\[
\frac{Fq}{A_0}=Y(q-1),
\qquad
F=A_0Y\left(1-\frac1q\right). \tag{23}
\]

这说明题目虽然把材料应力对工程应变写成线性形式，整根绳的力---伸长关系却因面积收缩而变成非线性。对
\(q\) 求导：

\[
\frac{dF}{dq}=\frac{A_0Y}{q^2}>0,
\qquad
\frac{d^2F}{dq^2}=-\frac{2A_0Y}{q^3}<0. \tag{24}
\]

拉力随伸长单调增大，但斜率逐渐减小，并从下方趋近
\(A_0Y\)。这一上限不是材料断裂强度，而是题设理想本构与体积约束共同产生的数学极限；真实材料在到达该极限之前通常早已超出线性弹性范围。

势能也可用无量纲形式写成

\[
E=A_0Yl_0(q-1-\ln q). \tag{25}
\]

函数 \(q-1-\ln q\ge0\)，等号只在 \(q=1\) 成立，因此势能非负。又

\[
\frac{dE}{dl}=F,
\]

可直接核验积分常数与拉力表达式彼此一致。

悬绳中真正连续的是材料标签而不是形变后空间坐标。用未形变坐标 \(\lambda\)
时，下方绳质量恰为
\(M\lambda/l_0\)；若改用形变后坐标，线密度本身随位置变化，必须额外求雅可比，反而更复杂。式（16）还揭示局部可行条件
\(T(\lambda)<A_0Y\)，只需检查张力最大的上端即可得到式（22）。

对小载荷作独立展开。令

\[
\epsilon_m=\frac{mg}{A_0Y},
\qquad
\epsilon_M=\frac{Mg}{A_0Y},
\]

且二者都远小于 \(1\)。式（21）中的对数为

\[
\ln\frac{1-\epsilon_m}{1-\epsilon_m-\epsilon_M}
\simeq\epsilon_M+\epsilon_m\epsilon_M+\frac12\epsilon_M^2. \tag{26}
\]

乘以前因子 \(l_0/\epsilon_M\)，得到

\[
l_2\simeq l_0\left[1+\frac{(m+M/2)g}{A_0Y}\right]. \tag{27}
\]

绳自重在一阶近似中等效为把 \(M/2\)
加到端点负载上，因为张力从下端到上端线性变化，其平均值包含半根绳的重量。这是式（21）非常有用的物理检查。

\needspace{3\baselineskip}
\textbf{材料坐标与图中的粗细}

令 \(F_*=A_0Y\)，从绳的下端沿原长量材料坐标
\(\lambda\)。截面张力和面积为

\[
T(\lambda)=mg+\frac{Mg}{l_0}\lambda,
\qquad A(\lambda)=A_0\left(1-\frac{T(\lambda)}{F_*}\right).
\]

面积随材料坐标线性减小，但空间高度不是
\(\lambda\)。从下端到该截面的实际弧长是

\[
s(\lambda)=\frac{F_*l_0}{Mg}
\ln\frac{F_*-mg}{F_*-mg-Mg\lambda/l_0}.
\]

对它求导，恰得
\(ds/d\lambda=F_*/[F_*-T(\lambda)]\)。图1-2用材料标签定位微元，上端画得更细，表示面积而非质量密度下降。可行条件是
\(F_*>(m+M)g\)；接近等号时，对数长度发散。这是题设本构模型的结果，不能当作真实绳子的断裂强度。

令 \(M\to0\) 时，应先对对数展开，得到 \(l_2\to l_0F_*/(F_*-mg)\)；不能把
\(M=0\) 直接代入含 \(1/M\) 的形式。

\needspace{3\baselineskip}
\textbf{一个可回代的悬绳样例与恒载稳定性}

为检查积分中的材料坐标，另设下端载荷占 \(F_*\) 的0.2、绳自重占 \(F_*\)
的0.3。记材料坐标比例 \(\xi=\lambda/l_0\)，局部伸长比和截面比分别为

\[
q(\xi)=\frac1{0.8-0.3\xi},\qquad
\frac{A(\xi)}{A_0}=0.8-0.3\xi.
\]

下端的q为1.25、上端为2；整个绳长不能简单取任一端的伸长比。正确积分为

\[
\frac l{l_0}=\int_0^1\frac{d\xi}{0.8-0.3\xi}
=\frac1{0.3}\ln\frac{0.8}{0.5}\simeq1.56668.
\]

它介于两端伸长比之间。把空间坐标解成材料标签后，任何截面的张力仍等于其下方真实质量的重量，不能对变细后的图形直接假定线密度恒定。图1-2中的粗细只用来提示面积变化，定量关系要由上述两式读取。

忽略绳自重时，还可检查恒外力P下的势能。令伸长状态势函数为

\[
\Pi(l)=F_*[l-l_0-l_0\ln(l/l_0)]-P(l-l_0).
\]

驻点条件 \(\Pi'=0\) 给 \(l=l_0F_*/(F_*-P)\)，而

\[
\Pi''=\frac{F_*l_0}{l^2}>0.
\]

在 \(P<F_*\)
下，这是题设模型的局部稳定伸长态；P接近F\emph{时长度趋于无穷，不是一个可代入有限绳长的普通端点。P不小于F}时，没有有限正长度平衡。这个判据来自本构式，不应与实际绳子的屈服、断裂判据混为一谈。

题目把杨氏关系与体积不变延伸到大形变，所得非线性力只是这套给定假设的形式结果。样例上端伸长一倍，已远超常规小应变测量区间，适合用于代数核验；若设计真实实验，要限制形变并提供相应材料的应力应变关系。讨论这一边界能帮助学生理解为何数学积分成立，材料适用性仍需另外检查。

\subsection{六、关键洞察}\label{ux516dux5173ux952eux6d1eux5bdf}

解析图区分材料坐标 \(\lambda\) 与伸长后的空间距离
\(s(\lambda)\)。悬点以下所含材料质量决定拉力，真实截面再由不可压缩条件求出。

\subsection{七、易错点}\label{ux4e03ux6613ux9519ux70b9}

\begin{enumerate}
\def\labelenumi{\arabic{enumi}.}
\tightlist
\item
  把 \(A\) 错当常量，写成 \(F=A_0Y(l-l_0)/l_0\)。
\item
  用 \(E=\tfrac12F_1(l_1-l_0)\)；这里 \(F(l)\) 不是线性函数。
\item
  把悬绳各处张力都写成 \((m+M)g\)。
\item
  积分时漏掉材料微元体积守恒，直接套用小形变公式。
\item
  忽略 \(A_0Y>(m+M)g\) 的存在条件。
\item
  把题设应变 \((l-l_0)/l_0\) 与真实应变 \(\ln(l/l_0)\)
  混用；本题必须严格使用给定本构关系。
\item
  在悬绳积分中用形变后长度直接分配质量，等价于错误地假设形变后线密度仍均匀。
\item
  对数积分上下限反写，得到负的绳长或负的伸长量。
\item
  看到 \(F\to A_0Y\)
  就误解为任何大长度都能在有限安全应力下实现，忽略模型适用范围。
\end{enumerate}

\subsection{八、课堂讲解建议}\label{ux516bux8bfeux5802ux8bb2ux89e3ux5efaux8bae}

先让学生比较``小形变细杆''和``体积保持的有限形变粗绳''两种模型，明确应力分母为何是当前面积。讲第二问时用图中的
\(\lambda\) 截面逐步增加下方绳段，帮助学生理解张力梯度，再进入积分。

板书可分成``整体均匀拉伸''和``非均匀悬挂''两栏。左栏从 \(Al=A_0l_0\)
出发，先求
\(F(l)\)，再画一条单调上升但逐渐变平的曲线，并用曲线下面积解释势能。右栏先画图
1-2，在任意截面处圈出下方系统，写出
\(T(\lambda)\)；随后只对一个材料微元推导
\(ds/d\lambda\)，最后积分整根绳。这样学生能看到第二问不是重新发明公式，而是把第一问的局部关系应用到每个微元。

讲积分换元时要求逐项写出上下限：\(\lambda=0\) 对应
\(A_0Y-mg\)，\(\lambda=l_0\) 对应 \(A_0Y-(m+M)g\)。因为新变量随
\(\lambda\)
减小，微分带负号；上下限与负号共同保证对数为正。可让学生先根据``悬挂后必有
\(l_2>l_0\)''预测结果符号，再做代数。

课堂验收至少包含三项回代：\(M\to0\) 时回到无自重端点负载；\(m,M\to0\) 时
\(l_2\to l_0\)；小载荷时回到式（27）。若学生只给精确对数式而没有任何极限检查，教师应追问其是否真正理解参数作用。

建议把学生常见的错误线性答案与精确答案放在同一图上比较。小伸长时两者斜率相同；伸长增大后，面积收缩使精确拉力曲线逐渐低于定面积直线。再让学生解释势能为什么不能写成末态力乘位移的一半：只有力---位移图严格为过原点直线时，三角形面积公式才成立。本题曲线下面积必须积分，图像与数学方法由此对应。

\subsection{九、变式训练}\label{ux4e5dux53d8ux5f0fux8badux7ec3}

若绳的自重可忽略，只悬挂质量 \(m\)，证明

\[
l=\frac{A_0Y}{A_0Y-mg}l_0.
\]

再对 \(mg\ll A_0Y\) 展开，说明它与线性胡克定律的一阶结果一致。

解析时令 \(M\to0\) 不能直接代入式（21），因为前因子和对数构成
\(0/0\)。可从局部伸长比直接得到

\[
q=\frac{A_0Y}{A_0Y-mg},
\]

也可对式（21）使用洛必达法则。\Needspace{20\baselineskip}

再令 \(x=mg/(A_0Y)\ll1\)，利用 \((1-x)^{-1}\simeq1+x\)，得到

\[
l-l_0\simeq\frac{mgl_0}{A_0Y},
\]

即等效劲度 \(A_0Y/l_0\)
的线性胡克定律结果。还可进一步训练：若上端和下端分别悬挂不同物体，先画受力再决定张力函数；若密度随伸长变化，体积约束应由新的质量守恒关系替代，不能沿用式（13）。

\newpage

\section{第 2 题}\label{ux7b2c-2-ux9898}

（1）半径为 \(a\) 的孤立金属小球放在匀强外电场 \(\boldsymbol E_0\)
中，求感应电荷形成的电偶极矩 \(\boldsymbol p_0\)。

（2）在相对介电常数 \(\varepsilon_1\simeq1\) 的介质中，均匀分布数密度为
\(n\)、半径均为 \(a\) 的金属小球，构成一种复合材料。已知

\[
\frac1n\sim\frac43\pi a^3,
\]

求复合材料的等效相对介电常数 \(\varepsilon_2\)。

\begin{figure}[H]
\centering
\includegraphics[width=0.76\linewidth,height=\textheight,keepaspectratio,alt={图2-1　导体球在外场中的极化及复合介质的局域场}]{../figures/question-02.pdf}
\caption*{图2-1　导体球在外场中的极化及复合介质的局域场}
\end{figure}

\subsection{一、题型定位}\label{ux4e00ux9898ux578bux5b9aux4f4d-1}

单个导体球的偶极矩可以由边界条件直接求出。由它推算复合介质时，还要说明局域场模型；题给的小球间距并不保证稀疏，平均场结果不能冒充任意排列的精确答案。

\subsection{二、知识点}\label{ux4e8cux77e5ux8bc6ux70b9-1}

\begin{enumerate}
\def\labelenumi{\arabic{enumi}.}
\tightlist
\item
  导体表面为等势面，球外电势可写成匀强场项与偶极项之和。
\item
  半径为 \(a\) 的导体球的静电极化率为 \(4\pi\varepsilon_0a^3\)。
\item
  各向同性平均场模型中的洛伦兹局域场为
  \(\boldsymbol E_{\rm loc}=\boldsymbol E+\boldsymbol P/(3\varepsilon_0)\)。
\item
  宏观极化强度满足 \(\boldsymbol P=n\boldsymbol p\)。
\item
  \(\boldsymbol P=\varepsilon_0(\varepsilon_2-1)\boldsymbol E\)。
\end{enumerate}

\subsection{三、物理图像}\label{ux4e09ux7269ux7406ux56feux50cf-1}

外场把导体球表面的自由电荷推向相反两侧；球内总电场为零，球外远处则恢复匀强场。复合材料中每个球不仅受到宏观平均场，还受到周围极化电荷产生的局域场。

\subsection{四、解题入口}\label{ux56dbux89e3ux9898ux5165ux53e3-1}

第一问从轴对称拉普拉斯方程的通解出发，用 \(r=a\)
为等势面确定偶极项系数。第二问把单球响应中的外场替换为局域场，建立关于
\(\boldsymbol P\) 的自洽方程。

\subsection{五、解析}\label{ux4e94ux89e3ux6790-1}

\begin{figure}[H]
\centering
\includegraphics[width=0.82\linewidth,height=\textheight,keepaspectratio,alt={图2-2　球外等势线与洛伦兹空腔场的自洽关系}]{../figures/solution-02.pdf}
\caption*{图2-2　球外等势线与洛伦兹空腔场的自洽关系}
\end{figure}

图2-2先画单球球外场，再画洛伦兹球形空腔。单球边界解是严格的，球腔场用于复合材料时另有平均场近似，不能由图形相似直接推出浓集材料的精确响应。

\needspace{3\baselineskip}
\textbf{（一）球外电势的一般形式}

取 \(\boldsymbol E_0\) 沿 \(z\)
轴。轴对称且在无电荷区域满足拉普拉斯方程的最低阶电势可写为

\[
\varphi_{\rm out}
=-E_0r\cos\theta+\frac{C\cos\theta}{r^2}. \tag{1}
\]

第一项是外加匀强场电势，第二项具有偶极电势的空间形式。

导体球表面必须为等势面。可把球面电势常数取为零，于是

\[
\varphi_{\rm out}(a,\theta)
=-E_0a\cos\theta+\frac{C}{a^2}\cos\theta=0. \tag{2}
\]

因为式（2）对任意 \(\theta\) 都成立，

\[
C=E_0a^3. \tag{3}
\]

所以

\[
\varphi_{\rm out}
=-E_0r\cos\theta
+E_0a^3\frac{\cos\theta}{r^2}. \tag{4}
\]

\needspace{3\baselineskip}
\textbf{（二）与偶极电势比较}

位于原点、方向沿 \(z\) 轴的电偶极子电势为

\[
\varphi_p
=\frac{1}{4\pi\varepsilon_0}
\frac{p\cos\theta}{r^2}. \tag{5}
\]

比较式（4）和式（5）的 \(r^{-2}\) 项：

\[
\frac{p_0}{4\pi\varepsilon_0}=E_0a^3. \tag{6}
\]

因此

\[
\boxed{\boldsymbol p_0=4\pi\varepsilon_0a^3\boldsymbol E_0}. \tag{7}
\]

偶极矩方向与外电场方向相同。

\needspace{3\baselineskip}
\textbf{（三）复合介质中的局域场}

在球形空腔平均场模型中，代表性小球的局域场近似取为

\[
\boldsymbol E_{\rm loc}
=\boldsymbol E+\frac{\boldsymbol P}{3\varepsilon_0}. \tag{8}
\]

单球偶极矩为

\[
\boldsymbol p
=4\pi\varepsilon_0a^3\boldsymbol E_{\rm loc}. \tag{9}
\]

单位体积内有 \(n\) 个球，所以

\[
\boldsymbol P=n\boldsymbol p. \tag{10}
\]

将式（8）、式（9）代入：

\[
\boldsymbol P
=4\pi n\varepsilon_0a^3
\left(
\boldsymbol E+\frac{\boldsymbol P}{3\varepsilon_0}
\right). \tag{11}
\]

定义

\[
f=\frac43\pi na^3, \tag{12}
\]

于是 \(4\pi na^3=3f\)，式（11）变为

\[
\boldsymbol P
=3f\varepsilon_0\boldsymbol E+f\boldsymbol P. \tag{13}
\]

移项：

\[
(1-f)\boldsymbol P
=3f\varepsilon_0\boldsymbol E. \tag{14}
\]

故

\[
\boldsymbol P
=\frac{3f}{1-f}\varepsilon_0\boldsymbol E. \tag{15}
\]

另一方面，等效介质定义给出

\[
\boldsymbol P
=\varepsilon_0(\varepsilon_2-1)\boldsymbol E. \tag{16}
\]

比较式（15）与式（16）：

\[
\varepsilon_2-1=\frac{3f}{1-f}. \tag{17}
\]

因此，在上述平均场近似下

\[
\boxed{
\varepsilon_2=\frac{1+2f}{1-f}
}. \tag{18}
\]

代回 \(f=\tfrac43\pi na^3\)：

\[
\boxed{
\varepsilon_2=
\frac{1+\dfrac83\pi na^3}
{1-\dfrac43\pi na^3}
}. \tag{19}
\]

若另取稀疏极限 \(f\ll1\)，则

\[
\varepsilon_2
=1+3f+O(f^2)
\simeq1+4\pi na^3. \tag{20}
\]

\needspace{3\baselineskip}
\textbf{补充推导与适用条件}

单球问题的关键是外场项与偶极项具有相同的角依赖
\(\cos\theta\)，所以只需用一个系数就能把整个球面变成等势面。球外径向电场为

\[
E_r=-\frac{\partial\varphi}{\partial r}
=E_0\cos\theta+2E_0a^3\frac{\cos\theta}{r^3}. \tag{21}
\]

在 \(r=a\) 处

\[
E_r(a)=3E_0\cos\theta. \tag{22}
\]

导体表面外侧法向场满足 \(\sigma=\varepsilon_0E_r\)，故

\[
\sigma(\theta)=3\varepsilon_0E_0\cos\theta. \tag{23}
\]

用表面电荷独立计算偶极矩的 \(z\) 分量：

\[
\begin{aligned}
p_z
&=\int z\,dq
=\int_0^{2\pi}\!\int_0^\pi
(a\cos\theta)\sigma(\theta)a^2\sin\theta\,d\theta d\varphi\\
&=6\pi\varepsilon_0E_0a^3
\int_0^\pi\cos^2\theta\sin\theta\,d\theta\\
&=4\pi\varepsilon_0a^3E_0.
\end{aligned} \tag{24}
\]

这与电势比较法完全一致，也验证了偶极矩方向。

复合介质中有三个不同的场：宏观平均场 \(\boldsymbol E\)、洛伦兹空腔局域场
\(\boldsymbol E_{\rm loc}\)、导体球内部的零场。式（8）连接前两个；导体球极化率把局域场映射为偶极矩；\(\boldsymbol P=n\boldsymbol p\)
再把微观响应带回宏观。这个闭环是自洽场方法的本质。

结果还可写成标准 Clausius--Mossotti 形式

\[
\frac{\varepsilon_2-1}{\varepsilon_2+2}=f. \tag{25}
\]

将式（18）代入即可验证。\(f\to0\) 时 \(\varepsilon_2\to1\)；小 \(f\) 下
\(\varepsilon_2-1\simeq3f\)，表明每增加少量导体体积分数都会增强介电响应。\(f\to1\)
时公式发散，但稀疏、互不接触的小球假设早已失效，不能把该发散解释为真实材料在任意结构下必然发生。

若基质相对介电常数为
\(\varepsilon_h\)，单球极化率应使用基质介电常数，导体极限为
\(4\pi\varepsilon_0\varepsilon_h a^3\)；宏观极化的定义也要相对基质背景重新组织。不能只在最终式中把
\(1\) 机械替换成 \(\varepsilon_h\) 而不改变局域场方程。

\needspace{3\baselineskip}
\textbf{题给浓度与平均场条件}

记 \(f=4\pi na^3/3\)。原题说 \(1/n\) 与一个球的体积同数量级，所以 \(f\)
未必很小。式（18）

\[
\varepsilon_2=\frac{1+2f}{1-f}
\]

是采用球形空腔局域场 \(E_{\rm loc}=E+P/(3\varepsilon_0)\)
后的平均场形式结果。单球偶极矩是严格解，浓集小球之间的多极相互作用却没有包含在这个闭合关系里。仅给
\(a,n\)，没有给球的排列、间隙和接触情况，不能唯一确定任意浓集材料的精确介电常数。

球形空腔贡献可直接算出：均匀极化沿 \(z\)，空腔内表面的束缚电荷为
\(\sigma_b=-P\cos\vartheta\)。其中心场为

\[
E_{{\rm cav},z}
=\frac{1}{4\pi\varepsilon_0}\int
\frac{-\sigma_b\cos\vartheta}{b^2}\,dS
=\frac{P}{4\pi\varepsilon_0}\int\cos^2\vartheta\,d\Omega
=\frac{P}{3\varepsilon_0}.
\]

这里算出的只是连续极化背景的空腔场；将其当作离散小球的全部局域修正，还需要忽略近邻结构场。图2-2的反馈箭头对应这个模型闭合步骤。另取
\(f\ll1\) 时，一阶结果 \(\varepsilon_2=1+3f+O(f^2)\)
与独立偶极子叠加相符；题给``同数量级''不能直接当成这个条件。若保留基质相对介电常数
\(\varepsilon_1\)，同一模型给
\(\varepsilon_{\rm eff}=\varepsilon_1(1+2f)/(1-f)\)，本题取
\(\varepsilon_1\simeq1\)。

\subsection{六、关键洞察}\label{ux516dux5173ux952eux6d1eux5bdf-1}

解析图先画单球球外场，再画洛伦兹球形空腔。单球边界解是严格的，球腔场用于复合材料时另有平均场近似，不能由图形相似直接推出浓集材料的精确响应。

\subsection{七、易错点}\label{ux4e03ux6613ux9519ux70b9-1}

\begin{enumerate}
\def\labelenumi{\arabic{enumi}.}
\tightlist
\item
  漏写 \(\varepsilon_0\)，导致偶极矩量纲错误。
\item
  把球外电势的偶极项错写成 \(1/r\) 或 \(1/r^3\)。
\item
  把 \(\boldsymbol P\) 与总偶极矩混淆；\(\boldsymbol P\)
  是单位体积偶极矩。
\item
  在自洽方程中漏掉 \(\boldsymbol P/(3\varepsilon_0)\)。
\item
  把近似式 \(1+3f\) 当成任意浓度下的精确式。
\item
  误把导体球内部场当作局域场；内部总场为零，局域场是球未放入前空腔中心的场。
\item
  积分表面偶极矩时漏掉位置因子 \(a\cos\theta\) 或面积元
  \(a^2\sin\theta,d\theta d\varphi\)。
\item
  把平均场公式在 \(f\to1\)
  附近的发散当作严格预言；高体积分数下球间关联与接触效应不能忽略。
\item
  基质不是真空时仍使用 \(4\pi\varepsilon_0a^3\)，漏掉
  \(\varepsilon_h\)。
\end{enumerate}

\subsection{八、课堂讲解建议}\label{ux516bux8bfeux5802ux8bb2ux89e3ux5efaux8bae-1}

可先通过``球面上电势必须与 \(\theta\) 无关''确定
\(C\)，让学生看到边界条件如何决定感应偶极矩。有效介质部分要明确宏观场、局域场和球内场是三个不同概念。

第一问建议采用``两条路线、一次互检''。主路线由拉普拉斯方程和等势边界求电势；互检路线由式（23）的表面电荷积分偶极矩。两条路线共用图
2-2 的极轴和角度约定，能把抽象的 \(r^{-2}\)
系数与正负表面电荷联系起来。讲解中应指出球面总感应电荷为零，因为
\(\int_0^\pi\cos\theta\sin\theta,d\theta=0\)。

第二问先暂不代 \(f\)，保留 \(4\pi na^3\)，让学生完成 \(\boldsymbol P\)
的移项；再用 \(4\pi na^3=3f\) 简化。这样分母 \(1-f\)
的来源清楚可见。若一开始直接写 Maxwell--Garnett
公式，学生往往无法区分它与简单体积加权平均。

课堂最后要求做三个检查：\(f=0\)
回到基质；一阶展开得到独立偶极子叠加；把式（18）代入式（25）左右相等。还可以画出
\(\varepsilon_2(f)\) 在稀疏区间的曲线，说明线性近似何时开始偏离。

教师可用``拿走中心小球，留下球形空腔''的想象解释洛伦兹场：宏观场之外，空腔表面束缚电荷再贡献
\(\boldsymbol P/(3\varepsilon_0)\)。这一结论依赖局部统计各向同性；若小球排列有强方向性或体积分数不再小，局域场会成为张量或受多体关联影响。讲清这个边界，学生才能理解
Maxwell--Garnett 公式为什么不是任意复合材料的精确普适定律。

\subsection{九、变式训练}\label{ux4e5dux53d8ux5f0fux8badux7ec3-1}

若基质的相对介电常数为 \(\varepsilon_h\) 而不近似为
\(1\)，查明并推导导体球极限下的 Maxwell--Garnett 关系

\[
\frac{\varepsilon_{\rm eff}-\varepsilon_h}
{\varepsilon_{\rm eff}+2\varepsilon_h}=f.
\]

解析时先写导体球在基质中的极化率

\[
\alpha_s=4\pi\varepsilon_0\varepsilon_h a^3,
\]

再用相对于基质的局域场建立自洽关系。由题给目标式直接解
\(\varepsilon_{\rm eff}\)：

\[
\varepsilon_{\rm eff}
=\varepsilon_h\frac{1+2f}{1-f}. \tag{26}
\]

令 \(\varepsilon_h=1\)
即回到式（18）。可继续追问：若夹杂物不是理想导体而是介电常数
\(\varepsilon_i\) 的小球，右端应变为

\[
f\frac{\varepsilon_i-\varepsilon_h}
{\varepsilon_i+2\varepsilon_h},
\]

导体极限 \(\varepsilon_i\to\infty\) 才趋于
\(f\)。训练目标是理解边界条件与极化率如何进入有效介质，而不是记忆一个分式。

\newpage

\section{第 3 题}\label{ux7b2c-3-ux9898}

刘老师在竖直方向运动，一颗质量为 \(dM\) 的花椒粒以相对地面速度 \(u\)
撞上并黏住他。取向上为正，在 \(t\) 时刻刘老师的质量、速度分别为
\(M,v\)，在 \(t+dt\) 时刻共同质量、速度分别为
\(M+dM,v+dv\)。刘老师受到合外力 \(F\)，花椒粒所受外力可忽略。

（1）证明密舍尔斯基方程

\[
F=M\frac{dv}{dt}+(v-u)\frac{dM}{dt}.
\]

（2）此后刘老师由静止下落，并团成密度为 \(\rho\)
的球体；他主动吸收空气中的水，质量增加率与球面面积成正比，比例系数为
\(k\)。初始半径为 \(r_0\)，求质量 \(M(t)\)。

（3）不计空气阻力，求自静止开始计时后的下落高度 \(h(t)\)。

（4）令
\(r_0\to0\)。升级后他不再主动吸收，只吸收运动过程中球体扫过体积内的水，水的密度为
\(\rho/2333\)。求升级后与升级前下落相同高度所需时间之比。

\begin{figure}[H]
\centering
\includegraphics[width=0.82\linewidth,height=\textheight,keepaspectratio,alt={图3-1　变质量碰撞与球体下落吸收质量的过程}]{../figures/question-03.pdf}
\caption*{图3-1　变质量碰撞与球体下落吸收质量的过程}
\end{figure}

\subsection{一、题型定位}\label{ux4e00ux9898ux578bux5b9aux4f4d-2}

吸收物质会同时改变质量和动量。先在地面系列动量收支，再分别求按时间增长和按下落距离增长的半径；零初始半径须通过正初始半径的极限理解。

\subsection{二、知识点}\label{ux4e8cux77e5ux8bc6ux70b9-2}

\begin{enumerate}
\def\labelenumi{\arabic{enumi}.}
\tightlist
\item
  变质量系统必须把流入质量在惯性系中的速度 \(u\) 写入动量收支。
\item
  球质量 \(M=\tfrac43\pi\rho r^3\)，球面面积 \(S=4\pi r^2\)。
\item
  吸收静止物质时，向下取正可写成 \(M\dot v=Mg-v\dot M\)。
\item
  也可写成更便于积分的形式 \(d(Mv)/dt=Mg\)。
\item
  初始尺度为零时，问题常出现幂律相似解。
\end{enumerate}

\subsection{三、物理图像}\label{ux4e09ux7269ux7406ux56feux50cf-2}

球体吸收原来静止的水后，必须把一部分原有动量分给新加入的质量，所以加速度小于
\(g\)。主动吸收模式中半径按时间线性增长；扫掠吸收模式中半径增长率与下落速度相关，因此半径与下落距离成正比。

\subsection{四、解题入口}\label{ux56dbux89e3ux9898ux5165ux53e3-2}

第一问直接对 \(t\) 到 \(t+dt\)
的``主体加上即将流入的微小质量''列动量定理。第二问先消去 \(M\) 得到
\(\dot r\)。第三问把密舍尔斯基方程改写为 \(d(Mv)/dt=Mg\)。第四问利用
\(r_0=0\) 的尺度齐次性设 \(h=At^2\)。

\subsection{五、解析}\label{ux4e94ux89e3ux6790-2}

\begin{figure}[H]
\centering
\includegraphics[width=0.88\linewidth,height=\textheight,keepaspectratio,alt={图3-2　主动吸收与扫掠吸收的质量流和动量流}]{../figures/solution-03.pdf}
\caption*{图3-2　主动吸收与扫掠吸收的质量流和动量流}
\end{figure}

图3-2对两种增长机制分别画出速度与下落距离。外界供水按时间增大半径，扫掠供水按路程增大半径，动量方程均需计入新水的零初始速度。

\needspace{3\baselineskip}
\textbf{（一）有限时间元内的动量收支}

在 \(t\) 时刻，把``刘老师''和即将在 \(dt\)
内并入的花椒粒合并作为研究系统。初动量为

\[
p_{\rm i}=Mv+u\,dM. \tag{1}
\]

\(dt\) 后花椒粒已经黏住，末动量为

\[
p_{\rm f}=(M+dM)(v+dv). \tag{2}
\]

展开并舍去二阶小量 \(dM\,dv\)：

\[
\begin{aligned}
p_{\rm f}
&=Mv+M\,dv+v\,dM+dM\,dv\\
&\simeq Mv+M\,dv+v\,dM.
\end{aligned} \tag{3}
\]

外力冲量为 \(Fdt\)，故

\[
Fdt=p_{\rm f}-p_{\rm i}. \tag{4}
\]

代入式（1）、式（3）：

\[
\begin{aligned}
Fdt
&=M\,dv+v\,dM-u\,dM\\
&=M\,dv+(v-u)dM.
\end{aligned} \tag{5}
\]

两边除以 \(dt\)：

\[
\boxed{
F=M\frac{dv}{dt}+(v-u)\frac{dM}{dt}
}. \tag{6}
\]

这里 \(u\)
必须是流入质量在同一惯性系中的速度，不能误写成它相对主体的速度。

\needspace{3\baselineskip}
\textbf{（二）主动吸收时的半径和质量}

题设给出

\[
\frac{dM}{dt}=kS=4\pi kr^2. \tag{7}
\]

球体密度恒为 \(\rho\)：

\[
M=\frac43\pi\rho r^3. \tag{8}
\]

对时间求导：

\[
\frac{dM}{dt}=4\pi\rho r^2\frac{dr}{dt}. \tag{9}
\]

比较式（7）和式（9）：

\[
4\pi\rho r^2\frac{dr}{dt}=4\pi kr^2. \tag{10}
\]

当 \(r>0\) 时约去 \(4\pi r^2\)：

\[
\frac{dr}{dt}=\frac{k}{\rho}. \tag{11}
\]

由 \(r(0)=r_0\)，

\[
r(t)=r_0+\frac{k}{\rho}t. \tag{12}
\]

于是

\[
\boxed{
M(t)=\frac43\pi\rho
\left(r_0+\frac{k}{\rho}t\right)^3
}. \tag{13}
\]

\needspace{3\baselineskip}
\textbf{（三）主动吸收模式的下落速度}

下落阶段改取竖直向下为正。被吸收的水在地面惯性系中原来静止，所以
\(u=0\)。式（6）化为

\[
Mg=M\frac{dv}{dt}+v\frac{dM}{dt}. \tag{14}
\]

右边正是乘积导数：

\[
\frac{d}{dt}(Mv)
=M\frac{dv}{dt}+v\frac{dM}{dt}. \tag{15}
\]

故

\[
\frac{d}{dt}(Mv)=Mg. \tag{16}
\]

从静止开始，\(v(0)=0\)，积分得

\[
M(t)v(t)=g\int_0^tM(\tau)\,d\tau. \tag{17}
\]

记

\[
b=\frac{k}{\rho},
\qquad
M=C(r_0+bt)^3,
\qquad
C=\frac43\pi\rho. \tag{18}
\]

则

\[
\begin{aligned}
\int_0^tM(\tau)\,d\tau
&=C\int_0^t(r_0+b\tau)^3d\tau\\
&=\frac{C}{4b}\left[(r_0+bt)^4-r_0^4\right].
\end{aligned} \tag{19}
\]

代回式（17），并约去 \(C\)：

\[
\boxed{
v(t)=
\frac{g}{4b}
\left[
(r_0+bt)-\frac{r_0^4}{(r_0+bt)^3}
\right]
}. \tag{20}
\]

\needspace{3\baselineskip}
\textbf{（四）对速度积分求高度}

下落高度

\[
h(t)=\int_0^t v(\tau)\,d\tau. \tag{21}
\]

使用 \(s=r_0+b\tau\)，\(ds=b\,d\tau\)：

\[
\begin{aligned}
h(t)
&=\frac{g}{4b^2}
\int_{r_0}^{r_0+bt}
\left(s-\frac{r_0^4}{s^3}\right)ds\\
&=\frac{g}{4b^2}
\left[
\frac{s^2}{2}+\frac{r_0^4}{2s^2}
\right]_{r_0}^{r_0+bt}.
\end{aligned} \tag{22}
\]

因此

\[
\boxed{
h(t)=\frac{g}{8b^2}
\left[
(r_0+bt)^2+\frac{r_0^4}{(r_0+bt)^2}-2r_0^2
\right],
\quad b=\frac{k}{\rho}
}. \tag{23}
\]

检查 \(t=0\)：方括号为零。对式（23）求导可回到式（20）。

当 \(r_0\to0\) 时，

\[
M\propto t^3,\qquad
v=\frac{gt}{4},\qquad
\boxed{h_{\rm old}=\frac{gt^2}{8}}. \tag{24}
\]

\needspace{3\baselineskip}
\textbf{（五）升级后的扫掠吸收}

水密度为

\[
\rho_w=\frac{\rho}{2333}. \tag{25}
\]

球体以速度 \(v\) 前进，单位时间扫过体积为 \(\pi r^2v\)，所以

\[
\frac{dM}{dt}
=\rho_w\pi r^2v
=\frac{\rho}{2333}\pi r^2v. \tag{26}
\]

另一方面仍有

\[
\frac{dM}{dt}
=4\pi\rho r^2\frac{dr}{dt}. \tag{27}
\]

比较两式：

\[
4\frac{dr}{dt}=\frac{v}{2333}. \tag{28}
\]

又 \(v=dh/dt\)，所以

\[
\frac{dr}{dh}=\frac{1}{4\times2333}. \tag{29}
\]

\(r(0)=0\) 给出

\[
r=\frac{h}{4\times2333},\qquad M\propto r^3\propto h^3. \tag{30}
\]

由于方程没有新的时间尺度，设

\[
h=At^2. \tag{31}
\]

则

\[
v=2At,\qquad M=C_hh^3=C_hA^3t^6. \tag{32}
\]

代入动量形式 \(d(Mv)/dt=Mg\)：

\[
\frac{d}{dt}\left(2C_hA^4t^7\right)
=C_hA^3t^6g. \tag{33}
\]

求导并约去 \(C_hA^3t^6\)：

\[
14A=g. \tag{34}
\]

所以

\[
\boxed{h_{\rm new}=\frac{gt^2}{14}}. \tag{35}
\]

同一下落高度 \(H\) 对应

\[
t_{\rm old}=\sqrt{\frac{8H}{g}},
\qquad
t_{\rm new}=\sqrt{\frac{14H}{g}}. \tag{36}
\]

因此

\[
\boxed{
\frac{t_{\rm new}}{t_{\rm old}}
=\sqrt{\frac{14}{8}}
=\frac{\sqrt7}{2}
}. \tag{37}
\]

\needspace{3\baselineskip}
\textbf{补充推导与适用条件}

密舍尔斯基方程的符号可用``为了带动新质量，主体要付出动量''来理解。取下落方向为正且流入物质原来静止，\(u=0\)，方程给

\[
M\dot v=Mg-v\dot M. \tag{38}
\]

质量增长时
\(\dot M>0\)，第二项总会减小加速度。若流入物质已以与主体同向的速度 \(u\)
运动，则阻滞项变成 \((v-u)\dot M\)；当 \(u=v\)
时，新质量无需再被加速，主体瞬时满足 \(\dot v=g\)；当 \(u>v\)
时，流入质量反而向主体输入额外动量。

主动吸收模式在 \(r_0\to0\) 时具有一个普适结论。若

\[
M(t)=Ct^p,
\]

则

\[
\frac{d}{dt}(Mv)=Mg
\]

从 \(0\) 积分到 \(t\) 给出

\[
Ct^pv=Cg\frac{t^{p+1}}{p+1},
\]

故

\[
v=\frac{gt}{p+1},
\qquad
h=\frac{gt^2}{2(p+1)}. \tag{39}
\]

本题 \(M\propto r^3\propto t^3\)，即 \(p=3\)，立即得到
\(v=gt/4\)、\(h=gt^2/8\)。这条一般公式可以独立核对式（24）。

扫掠吸收模式则由 \(r\propto h\) 得 \(M\propto h^q\)，本题 \(q=3\)。设
\(h=At^2\)，则

\[
M=C_hA^qt^{2q},
\qquad
Mv=2C_hA^{q+1}t^{2q+1}. \tag{40}
\]

代入动量方程并约去公共因子：

\[
2A(2q+1)=g.
\]

所以一般结果为

\[
h=\frac{gt^2}{2(2q+1)}. \tag{41}
\]

取 \(q=3\) 就是 \(gt^2/14\)。这也清楚说明水密度只进入 \(C_h\)，而
\(C_h\) 在两边消去，时间比与 \(2333\) 的具体数值无关。数值 \(2333\)
只决定同一高度下球体究竟有多大、质量有多大。

式（23）还可在短时间展开。\(t\to0\) 时 \(M\simeq M_0\)，吸收量很小，应有
\(h\simeq gt^2/2\)。将 \(r_0+bt\) 展开到最低非零阶，式（20）确实给
\(v\simeq gt\)，式（23）给 \(h\simeq gt^2/2\)。长时间 \(bt\gg r_0\)
时则过渡到
\(h\simeq gt^2/8\)。这个早期---晚期极限能检查复杂分式中的系数和符号。

\needspace{3\baselineskip}
\textbf{用正初始半径检验零半径极限}

主动吸收时记 \(b=k/\rho\)，半径为 \(r_0+bt\)。已有位移式在
\(r_0\to0^+\)、固定 \(t>0\) 时趋于
\(gt^2/8\)，这是从有质量的初态取出的极限。

扫掠模式取 \(\kappa=\rho_w/(4\rho)=1/(4\times2333)\)，于是
\(r=r_0+\kappa h\)。对任意 \(r_0>0\)，由动量方程和 \(w(h)=v^2\) 得

\[
\frac{dw}{dh}+\frac{6\kappa}{r_0+\kappa h}w=2g,
\qquad w(0)=0.
\]

积分因子为 \((r_0+\kappa h)^6\)，所以

\[
v^2(h)=\frac{2g}{7\kappa}
\left[r_0+\kappa h-\frac{r_0^7}{(r_0+\kappa h)^6}\right].
\]

初始很短距离内 \(v^2\simeq2gh\)，因为新吸收质量尚少。取 \(r_0\to0^+\)
后，固定正高度有 \(v^2\to2gh/7\)；积分
\(dh/dt=\sqrt{2gh/7}\)，得到从释放时刻起的极限轨迹
\(h=gt^2/14\)。因此同一正高度的时间比为 \(\sqrt7/2\)。

严格令 \(r_0=0\)
会使起始质量与截面积同时为零，常规初值方程奇异；原题结论应理解为上述正种子极限，不应宣称零质量初态有唯一普通解。图3-2的两条抛物线据此画出，并标明是
\(r_0\to0^+\)
的轨迹。水在地面系静止是另一个必要假设；若水已有速度，要恢复
\((v-u)\dot M\) 项。

\subsection{六、关键洞察}\label{ux516dux5173ux952eux6d1eux5bdf-2}

解析图对两种增长机制分别画出速度与下落距离。外界供水按时间增大半径，扫掠供水按路程增大半径，动量方程均需计入新水的零初始速度。

\subsection{七、易错点}\label{ux4e03ux6613ux9519ux70b9-2}

\begin{enumerate}
\def\labelenumi{\arabic{enumi}.}
\tightlist
\item
  把流入质量速度 \(u\) 写成相对速度。
\item
  展开 \((M+dM)(v+dv)\) 时保留二阶小量，或漏掉 \(v\,dM\)。
\item
  认为吸收静止水不影响运动而直接令 \(a=g\)。
\item
  把扫过横截面积误写成球面面积 \(4\pi r^2\)。
\item
  用主动吸收的 \(r\propto t\) 结论套到升级模式。
\item
  在同一推导中先取向上为正、后又把 \(g\)
  当正数代入，没有明确切换坐标正向。
\item
  把 \(d(Mv)/dt=Mg\) 积分成 \(Mv=Mgt\)，忽略 \(M\)
  本身随时间变化，正确右端是 \(g\int Mdt\)。
\item
  扫掠吸收中使用球面面积 \(4\pi r^2\)；扫过体积由迎风截面 \(\pi r^2\)
  决定。
\item
  把严格零质量初态当作普通初值问题，或只写一个相似解便宣称它唯一；应从
  \(r_0>0\) 的解取极限。
\end{enumerate}

\subsection{八、课堂讲解建议}\label{ux516bux8bfeux5802ux8bb2ux89e3ux5efaux8bae-2}

先用一维小车接沙模型解释 \(-v\dot M\)
项的物理来源。第四问不要一上来做繁琐联立微分方程，可先问学生``若
\(r\propto h\)，质量随高度是几次方''，再由量纲和幂律试探建立相似解。

第一问板书应明确研究系统：在 \(t\) 时刻把主体与即将在 \(dt\)
内并入的质量一起选入系统。初态中两部分速度不同，末态中速度相同；这样
\(u,dM\) 的来源一目了然。展开末动量时用括号标出二阶小量
\(dM,dv\)，说明其相对 \(dt\) 为高阶，才能安全舍去。

第二、三问建议沿``几何增长---质量函数---动量积分---位移积分''的顺序。得到式（20）后先做
\(t=0\) 和长时间极限，再积分位移。积分中的 \(r_0^4/s^3\)
容易写错符号，可由对最终 \(h(t)\) 求导回到 \(v(t)\) 检查。

第四问先在图 3-2
上区分``主动从整个表面吸收''和``运动扫过迎风截面吸收''。由两种 \(dM/dt\)
表达式相等得到 \(dr/dh\)，这是最关键的一步。随后只保留幂次关系
\(M\propto h^3\)，再做相似解，避免让常数 \(\rho,2333,\pi\)
遮蔽物理结构。

课堂验收应要求学生解释为什么升级后反而下落更慢：升级模式中随高度增长，质量按
\(h^3\)
快速增大，主体必须不断把静止水加速到自身速度；被吸收物质的动量负担使
\(h\) 的系数从 \(1/8\) 降到 \(1/14\)。

还应让学生对两种模式分别画 \(M(t)\)、\(v(t)\) 的幂律图。主动模式长时间
\(M\propto t^3\)、\(v\propto t\)；扫掠模式
\(h\propto t^2\)、\(M\propto h^3\propto t^6\)、\(v\propto t\)。速度同为时间一次函数，但斜率不同，原因是质量增长幂次不同。把这些幂次代回
\(d(Mv)/dt=Mg\)，左右都应为同次幂，是相似解成立的必要检查。

\subsection{九、变式训练}\label{ux4e5dux53d8ux5f0fux8badux7ec3-2}

若升级模式吸收的是随高度方向以恒定速度 \(u_0\)
下落的水，写出相应的密舍尔斯基方程，并讨论 \(u_0\)
与球速同向或反向时质量流对加速度的影响。

取下落为正，流入水的地面速度为 \(u_0\)，则

\[
Mg=M\dot v+(v-u_0)\dot M,
\]

或

\[
\frac{d}{dt}(Mv)=Mg+u_0\dot M. \tag{42}
\]

若 \(0<u_0<v\)，阻滞仍存在但弱于静水；若
\(u_0=v\)，质量加入时无相对速度，瞬时不产生变质量阻滞；若
\(u_0>v\)，流入水追上球体并输入动量，质量流项可提高加速度。若
\(u_0<0\)，水向上运动，主体还必须反转其动量，阻滞最强。可进一步让学生讨论
\(u_0\) 为常量时原来的 \(h=At^2\)
纯相似形式会被新的速度尺度破坏，通常需要解含平移项的微分方程。

\newpage

\section{第 4 题}\label{ux7b2c-4-ux9898}

半径为 \(r\)、匝数为 \(n\)
的线圈无摩擦地套在磁场方向呈辐射状的永磁体柱上，线圈处磁感应强度为
\(B\)。线圈与小灯泡构成总电阻为 \(R\)
的闭合回路。三个相同轻弹簧的劲度系数均为
\(k/3\)，线圈及套筒质量忽略。平衡时弹簧均处于原长，记线圈在平衡位置以下的位移为
\(x\)。

（1）施加恒定外力 \(F_0\)，使线圈从 \(x=F_0/(3k)\) 运动到
\(x=2F_0/(3k)\)；随后撤去外力，待线圈回到
\(x=F_0/(3k)\)。求一个循环内小灯泡的平均功率 \(P_0\)。

（2）改变外力，使线圈以恒定速率 \(v_0\) 从 \(F_0/(3k)\) 向下运动到
\(2F_0/(3k)\)，再以相同速率向上返回。若平均功率仍为 \(P_0\)，求
\(v_0\)。

\begin{figure}[H]
\centering
\includegraphics[width=0.632\linewidth,height=\textheight,keepaspectratio,alt={图4-1　辐射磁场中的线圈、灯泡和三弹簧发电装置}]{../figures/question-04.pdf}
\caption*{图4-1　辐射磁场中的线圈、灯泡和三弹簧发电装置}
\end{figure}

\subsection{一、题型定位}\label{ux4e00ux9898ux578bux5b9aux4f4d-3}

辐射磁场使线圈产生与速度成正比的电磁阻力。质量和自感忽略后，速度由瞬时力平衡决定；循环功、循环时间和灯泡所占电阻比例要分别核对。

\subsection{二、知识点}\label{ux4e8cux77e5ux8bc6ux70b9-3}

\begin{enumerate}
\def\labelenumi{\arabic{enumi}.}
\tightlist
\item
  线圈有效导线总长度为 \(2\pi rn\)。
\item
  动生电动势 \(\mathcal E=2\pi rnBv\)。
\item
  安培阻力与速度成正比，\(F_{\rm em}=(2\pi rnB)^2v/R\)。
\item
  三个并联弹簧的总劲度为 \(k\)。
\item
  质量忽略时没有惯性项，任意时刻满足代数力平衡。
\item
  一个完整循环中系统机械能回到原值，外力做功最终都转化为焦耳热。
\end{enumerate}

\subsection{三、物理图像}\label{ux4e09ux7269ux7406ux56feux50cf-3}

线圈上下运动切割辐射磁场，回路中产生电流；楞次定律使安培力总与速度反向。由于运动部件质量忽略，撤去外力后线圈不是自由振动，而是按一阶指数规律缓慢回到平衡位置。

\subsection{四、解题入口}\label{ux56dbux89e3ux9898ux5165ux53e3-3}

先把磁---电---力三个关系压缩为电磁阻尼系数
\(\gamma=(2\pi rnB)^2/R\)。恒力阶段由 \(kx+\gamma v=F_0\)
求时间，撤力阶段由 \(kx+\gamma v=0\)
求时间。平均功率可直接用``循环外力功除以周期''。

\subsection{五、解析}\label{ux4e94ux89e3ux6790-3}

\begin{figure}[H]
\centering
\includegraphics[width=0.88\linewidth,height=\textheight,keepaspectratio,alt={图4-2　恒力下行与撤力回程的一阶相轨迹}]{../figures/solution-04.pdf}
\caption*{图4-2　恒力下行与撤力回程的一阶相轨迹}
\end{figure}

图4-2只画 \(x=F_0/(3k)\) 到 \(2F_0/(3k)\)
的实际行程。下、上行两条速度线端点不同，积分时间却相同；图的有向循环对应输入功。

\needspace{3\baselineskip}
\textbf{（一）动生电动势、电流和安培阻力}

定义

\[
c=2\pi rnB. \tag{1}
\]

线圈以速度 \(v\) 沿轴向运动时，每匝有效导线长度为 \(2\pi r\)，故

\[
\mathcal E=cv. \tag{2}
\]

回路电流大小

\[
I=\frac{\mathcal E}{R}=\frac{cv}{R}. \tag{3}
\]

安培力大小等于总有效长度乘 \(IB\)：

\[
F_{\rm em}
=2\pi rn\,IB
=cI
=\frac{c^2}{R}v. \tag{4}
\]

令电磁阻尼系数

\[
\gamma=\frac{c^2}{R}
=\frac{(2\pi rnB)^2}{R}. \tag{5}
\]

则 \(F_{\rm em}=\gamma v\)，方向始终与运动方向相反。

\needspace{3\baselineskip}
\textbf{（二）恒力下行阶段}

取向下为正。三个弹簧总恢复力为 \(kx\)。质量忽略意味着

\[
F_0-kx-\gamma v=0. \tag{6}
\]

所以

\[
v=\frac{dx}{dt}=\frac{F_0-kx}{\gamma}. \tag{7}
\]

下行时间

\[
t_\downarrow
=\int_{F_0/(3k)}^{2F_0/(3k)}
\frac{\gamma\,dx}{F_0-kx}. \tag{8}
\]

令 \(u=F_0-kx\)，\(du=-kdx\)：

\[
\begin{aligned}
t_\downarrow
&=-\frac{\gamma}{k}
\int_{2F_0/3}^{F_0/3}\frac{du}{u}\\
&=\frac{\gamma}{k}\ln2.
\end{aligned} \tag{9}
\]

\needspace{3\baselineskip}
\textbf{（三）撤力回程阶段}

撤去外力后，线圈向上运动。若仍取向下为正，则 \(v<0\)。力平衡为

\[
-kx-\gamma v=0. \tag{10}
\]

因此

\[
v=-\frac{kx}{\gamma}. \tag{11}
\]

从 \(x=2F_0/(3k)\) 回到 \(x=F_0/(3k)\) 的时间：

\[
\begin{aligned}
t_\uparrow
&=\int_{2F_0/(3k)}^{F_0/(3k)}
\frac{dx}{-kx/\gamma}\\
&=-\frac{\gamma}{k}
\int_{2F_0/(3k)}^{F_0/(3k)}\frac{dx}{x}\\
&=\frac{\gamma}{k}\ln2.
\end{aligned} \tag{12}
\]

所以周期

\[
T=t_\downarrow+t_\uparrow
=\frac{2\gamma}{k}\ln2. \tag{13}
\]

\needspace{3\baselineskip}
\textbf{（四）用能量收支求平均功率}

完整循环前后线圈位置相同，弹性势能不变；质量忽略，也没有动能净变化。因此外力做功全部转化为灯泡和回路中的焦耳热。

恒力只在下行阶段做功，位移为

\[
\Delta x
=\frac{2F_0}{3k}-\frac{F_0}{3k}
=\frac{F_0}{3k}. \tag{14}
\]

每周期外力功

\[
W=F_0\Delta x=\frac{F_0^2}{3k}. \tag{15}
\]

平均功率

\[
\begin{aligned}
P_0
&=\frac{W}{T}\\
&=\frac{F_0^2/(3k)}
{(2\gamma/k)\ln2}\\
&=\frac{F_0^2}{6\gamma\ln2}.
\end{aligned} \tag{16}
\]

代入 \(\gamma=c^2/R\)：

\[
\boxed{
P_0=
\frac{F_0^2R}
{6\ln2\,(2\pi rnB)^2}
}. \tag{17}
\]

\needspace{3\baselineskip}
\textbf{（五）匀速往返模式}

速率恒为 \(v_0\) 时，电流大小在下行、上行阶段相同：

\[
I_0=\frac{cv_0}{R}. \tag{18}
\]

灯泡和回路的焦耳功率为常量

\[
P=I_0^2R
=\frac{c^2v_0^2}{R}. \tag{19}
\]

要求 \(P=P_0\)：

\[
\frac{c^2v_0^2}{R}
=\frac{F_0^2R}{6\ln2\,c^2}. \tag{20}
\]

解得

\[
\boxed{
v_0=
\frac{F_0R}
{\sqrt{6\ln2}\,(2\pi rnB)^2}
}. \tag{21}
\]

\needspace{3\baselineskip}
\textbf{补充推导与适用条件}

定义机电耦合常数 \(c=2\pi rnB\) 后，整个装置可压缩成

\[
\mathcal E=cv,
\qquad
I=\frac{cv}{R},
\qquad
F_{\rm em}=\frac{c^2}{R}v. \tag{22}
\]

三式满足

\[
F_{\rm em}v
=\frac{c^2v^2}{R}
=\mathcal EI
=I^2R, \tag{23}
\]

即机械功率、感应电源功率和焦耳热严格相等。这是阻尼系数中 \(c^2/R\)
而不是 \(c/R\) 的独立检查。

恒力阶段的速度

\[
v_\downarrow(x)=\frac{F_0-kx}{\gamma}
\]

在本题区间内从 \(2F_0/(3\gamma)\) 降到 \(F_0/(3\gamma)\)；回程速度大小

\[
|v_\uparrow(x)|=\frac{kx}{\gamma}
\]

从 \(2F_0/(3\gamma)\) 降到
\(F_0/(3\gamma)\)。两阶段沿相反方向经过同一位置区间，速度大小的端点次序互换，因此两个时间都等于
\((\gamma/k)\ln2\)。这不是偶然的``来回路程相等''，而是两个一阶分母分别从
\(2F_0/3\) 变到 \(F_0/3\)。

还可直接积分焦耳热验证每周期能量。下行阶段

\[
\begin{aligned}
Q_\downarrow
&=\int I^2R\,dt
=\int\gamma v^2\frac{dx}{v}\\
&=\int_{F_0/(3k)}^{2F_0/(3k)}
(F_0-kx)\,dx
=\frac{F_0^2}{6k}.
\end{aligned} \tag{24}
\]

回程取实际速度与位移同号，耗散仍为正：

\[
Q_\uparrow
=\int_{2F_0/(3k)}^{F_0/(3k)}(-kx)\,dx
=\frac{F_0^2}{6k}. \tag{25}
\]

两者相加
\(Q=F_0^2/(3k)\)，正好等于式（15）的外力功。这个逐段积分既核验能量法，也说明回程期间储存的弹性势能转化为焦耳热。

``质量忽略''应理解为机械时间常数远小于题目关注时间尺度。在真实装置中若线圈质量为
\(m_c\)，方程为

\[
m_c\ddot x+\gamma\dot x+kx=F(t). \tag{26}
\]

只有惯性项相对阻尼与弹性项可忽略时，才退化为本题一阶方程。不能把质量取零后仍假想线圈具有独立动能或越过平衡点振荡。

\needspace{3\baselineskip}
\textbf{灯泡功率与理想化条件}

设灯泡电阻为 \(R_b\)，其余导线和线圈电阻为
\(R-R_b\)。同一电流通过整个回路，所以

\[
\overline P_b=\frac{R_b}{R}\overline P_{\rm tot}
=\frac{R_b}{R}\frac{F_0^2R}{6\ln2\,(2\pi rnB)^2}.
\]

参考解把总焦耳热当作灯泡耗能，相当于认为线圈和导线电阻可忽略、\(R_b=R\)。若只知道总电阻而不知道分配，能确定的是总平均功率；灯泡功率须乘以上面的比例。两种工作方式使用同一回路时，这个比例在比较中约掉，所以匹配平均功率所求的
\(v_0\) 不变。

另一个默认条件是忽略线圈自感。若记自感为 \(L_s\)，电路应满足
\(L_s\dot I+RI=(2\pi rnB)\dot x\)。代数关系 \(I=(2\pi rnB)\dot x/R\)
要求电气暂态足够短。记电磁阻尼
\(\Gamma=(2\pi rnB)^2/R\)，题解使用的时间尺度为 \(\Gamma/k\)；可用
\(L_s/R\ll\Gamma/k\) 和 \(m_c/\Gamma\ll\Gamma/k\)
表示相应的近似条件。切换外力时理想模型允许速度立即改变，真实装置则有短暂过渡。

匀速往返时，向下坐标中的外力分别为 \(F_{\downarrow}=kx+\Gamma v_0\) 与
\(F_{\uparrow}=kx-\Gamma v_0\)，后式随位置可能为负。两条力曲线的间距为
\(2\Gamma v_0\)，一个循环外力功是 \(2\Gamma v_0\Delta x\)，时间为
\(2\Delta x/v_0\)，再得总平均功率
\(\Gamma v_0^2\)。图4-2的面积和两段指数运动分别对应这两种计算。

\subsection{六、关键洞察}\label{ux516dux5173ux952eux6d1eux5bdf-3}

解析图只画 \(x=F_0/(3k)\) 到 \(2F_0/(3k)\)
的实际行程。下、上行两条速度线端点不同，积分时间却相同；图的有向循环对应输入功。

\subsection{七、易错点}\label{ux4e03ux6613ux9519ux70b9-3}

\begin{enumerate}
\def\labelenumi{\arabic{enumi}.}
\tightlist
\item
  忘记 \(n\) 匝，或把有效长度只写成 \(2\pi r\)。
\item
  漏掉安培力中的第二个 \(B\)，使阻尼系数少一个 \(B\)。
\item
  把三个 \(k/3\) 弹簧串联；它们共同托起线圈，应按并联相加为 \(k\)。
\item
  撤力回程时符号处理错误。
\item
  只计算下行时间，漏掉回程的 \(\ln2\)。
\item
  把总电阻 \(R\)
  放到阻尼系数分子，忘记电阻越大、感应电流越小、阻尼越弱。
\item
  以为往返匀速时平均功率还需除以两个行程时间；瞬时焦耳功率已经恒定，平均值就是它本身。
\item
  撤去外力后仍把 \(F_0\)
  留在力平衡式中，或把向上速度当成正数却不改变坐标约定。
\item
  使用循环能量守恒求热量后，忘记还必须由动力学求周期才能得到平均功率。
\end{enumerate}

\subsection{八、课堂讲解建议}\label{ux516bux8bfeux5802ux8bb2ux89e3ux5efaux8bae-3}

先用功率链

\[
F_{\rm em}v=\mathcal EI=I^2R
\]

核验各式，帮助学生把机械功率、电磁功率和焦耳热统一起来。随后强调``质量可忽略''在动力学上的真正含义。

图 4-2 应与方程同步使用：在下行曲线上标出 \(x=F_0/(3k)\) 与
\(2F_0/(3k)\)
的速度，在回程曲线上标出负速度。让学生从图上先判断两阶段都越接近各自平衡位置越慢，再做分离变量积分。这样
\(\ln2\) 的出现有清楚的速度---位置来源。

求平均功率时可先让学生尝试逐时积分，再指出完整循环的初末弹性势能相同，外力只在下行做功，能量法立即给出总热量。随后用式（24）、式（25）作教师版互检。两种路线各有用途：能量法快，逐段积分能检查符号与能量去向。

第二问要强调匀速运动需要外力随位置和方向连续调节，但题目只问电功率，不必求该外力。因为
\(|v|\) 恒定，\(|\mathcal E|\)、\(I^2R\)
也恒定，方向反转只使电流反号，不改变焦耳功率。

验收时要求学生检查量纲：\(\gamma=c^2/R\) 应为
\(\mathrm{N\,s/m}\)，\(F_0/\gamma\) 是速度，\(F_0^2/\gamma\)
是功率。若最终 \(v_0\) 随 \(R\)
减小而减小，也符合强电流、强阻尼时达到同一功率只需较低速度的趋势。

可以设置一个口头追问：``电阻增大，为什么式（17）的 \(P_0\)
反而增大？''答案是第一问并非固定运动周期；电阻增大使电磁阻尼减弱，往返更快，虽然同一速度下功率减小，但每周期固定外力功在更短时间内耗散，平均功率反而增大。这个讨论能训练学生区分控制条件，避免只凭
\(P=\mathcal E^2/R\) 对最终趋势下结论。

\subsection{九、变式训练}\label{ux4e5dux53d8ux5f0fux8badux7ec3-3}

若线圈与套筒总质量为
\(m_c\)，写出恒力阶段的二阶微分方程，并判断系统在什么参数条件下会出现振荡回程。

恒力阶段的方程为式（26），平衡位置 \(x_e=F_0/k\)。令 \(y=x-x_e\)，得到

\[
m_c\ddot y+\gamma\dot y+ky=0. \tag{27}
\]

特征方程

\[
m_cr^2+\gamma r+k=0
\]

的判别式为 \(\gamma^2-4m_ck\)。当

\[
\gamma^2<4m_ck
\]

时为欠阻尼，线圈会绕平衡位置振荡；等号对应临界阻尼；大于零为过阻尼，不发生越过平衡位置的振荡。撤力后只需把平衡位置改为零，齐次参数不变。进一步可让学生比较
\(m_c\to0\)
时两个特征根：一个快速衰减模趋于无穷快，剩下的慢模正是本题的一阶解。

\newpage

\section{第 5 题}\label{ux7b2c-5-ux9898}

质量分别为 \(m_1,m_2\)
的两个原子组成双原子分子，只考虑连线方向的一维振动。两原子坐标为
\(x_1,x_2\)，间距 \(x=x_1-x_2>0\)。

（1）在质心系中，给出与相对加速度 \(\ddot x\) 对应的折合质量
\(\mu\)，并说明其来源。

（2）若平衡间距为 \(R\)，要求
\(x=R+C_1\cos\omega t+C_2\sin\omega t\)，写出 \(R\) 邻域内的势能增量
\(\Delta U(x)\)。

（3）实际势能近似为 Morse 势

\[
U(x)=D\left[e^{-2\alpha(x-R)}-2e^{-\alpha(x-R)}\right],
\]

其中 \(D=|U_{\min}|\)。机械能写成 \(E=-D+\varepsilon\)，求分子键不断裂时
\(\varepsilon\) 的范围。

（4）求平衡位置附近小振动角频率，并用 \(\varepsilon,\alpha,D\)
表示振幅。

（5）对一氧化碳，\(m_{\rm C}=12.00u\)，\(m_{\rm O}=15.99u\)，\(u=1.660\times10^{-27}\,\mathrm{kg}\)；测得
\(\omega=4.088\times10^{14}\,\mathrm{rad/s}\)，\(D=1.780\times10^{-18}\,\mathrm J\)，\(R=1.128\times10^{-10}\,\mathrm m\)。求
\(\alpha\) 及吸收波长，并判断波段。

思考题（不计分）：把分子中的一个原子替换为它的同位素会改变振动特征。比较
\(\mathrm{CH_4}\) 中的氕氢与 \(\mathrm{CH_3D}\)
中的氘氢参与卤化反应时的难易，判断哪一种键更容易发生反应。

\begin{figure}[H]
\centering
\includegraphics[width=0.8\linewidth,height=\textheight,keepaspectratio,alt={图5-1　双原子分子的相对坐标与 Morse 势阱}]{../figures/question-05.pdf}
\caption*{图5-1　双原子分子的相对坐标与 Morse 势阱}
\end{figure}

\subsection{一、题型定位}\label{ux4e00ux9898ux578bux5b9aux4f4d-4}

用两原子的相对距离作为坐标，惯性系数就是折合质量。Morse
势的曲率给小振动频率，远处的势能极限给解离条件，两者适用的能量范围不同。

\subsection{二、知识点}\label{ux4e8cux77e5ux8bc6ux70b9-4}

\begin{enumerate}
\def\labelenumi{\arabic{enumi}.}
\tightlist
\item
  两体内力问题可用质心坐标和相对坐标分离。
\item
  折合质量 \(\mu=m_1m_2/(m_1+m_2)\)。
\item
  稳定平衡点附近势能的二阶项决定小振动频率。
\item
  Morse 势在无穷远趋于零，解离阈值为 \(E=0\)。
\item
  角频率与真空波长关系 \(\lambda=2\pi c/\omega\)。
\end{enumerate}

\subsection{三、物理图像}\label{ux4e09ux7269ux7406ux56feux50cf-4}

Morse
势左侧陡峭、右侧平缓，体现原子过近时强排斥、过远时逐渐解离。小振幅时只看到势阱底部的抛物线，所以近似简谐；能量接近零时振动强烈非线性，最终分子键断裂。

\subsection{四、解题入口}\label{ux56dbux89e3ux9898ux5165ux53e3-4}

先在质心系把 \(x_1,x_2\) 写成相对坐标 \(x\)
的倍数，求动能中的等效质量。再令 \(q=x-R\)，对 Morse 势展开到
\(q^2\)。数值部分由 \(\omega=\alpha\sqrt{2D/\mu}\) 反求 \(\alpha\)。

\subsection{五、解析}\label{ux4e94ux89e3ux6790-4}

\begin{figure}[H]
\centering
\includegraphics[width=0.9\linewidth,height=\textheight,keepaspectratio,alt={图5-2　Morse 势的解离阈值、谐振近似和振幅定义}]{../figures/solution-05.pdf}
\caption*{图5-2　Morse 势的解离阈值、谐振近似和振幅定义}
\end{figure}

图5-2中Morse曲线与谐振抛物线只在谷底接近；有限能量水平线的左右交点不对称。后面的转折点、周期与能量关系据此计算。

\needspace{3\baselineskip}
\textbf{（一）折合质量}

质心系满足

\[
m_1x_1+m_2x_2=0. \tag{1}
\]

又

\[
x=x_1-x_2. \tag{2}
\]

联立解得

\[
x_1=\frac{m_2}{m_1+m_2}x,
\qquad
x_2=-\frac{m_1}{m_1+m_2}x. \tag{3}
\]

求导：

\[
\dot x_1=\frac{m_2}{m_1+m_2}\dot x,
\qquad
\dot x_2=-\frac{m_1}{m_1+m_2}\dot x. \tag{4}
\]

总动能

\[
\begin{aligned}
T
&=\frac12m_1\dot x_1^2+\frac12m_2\dot x_2^2\\
&=\frac12
\left[
\frac{m_1m_2^2+m_2m_1^2}{(m_1+m_2)^2}
\right]\dot x^2\\
&=\frac12
\frac{m_1m_2(m_1+m_2)}{(m_1+m_2)^2}
\dot x^2\\
&=\frac12\mu\dot x^2,
\end{aligned} \tag{5}
\]

其中

\[
\boxed{\mu=\frac{m_1m_2}{m_1+m_2}}. \tag{6}
\]

于是相对坐标的运动方程为

\[
\mu\ddot x=-\frac{dU}{dx}. \tag{7}
\]

\needspace{3\baselineskip}
\textbf{（二）严格简谐运动对应的势能}

题给

\[
x-R=C_1\cos\omega t+C_2\sin\omega t. \tag{8}
\]

两次求导：

\[
\ddot x=-\omega^2(x-R). \tag{9}
\]

代入式（7）：

\[
-\frac{dU}{dx}=-\mu\omega^2(x-R). \tag{10}
\]

故

\[
\frac{dU}{dx}=\mu\omega^2(x-R). \tag{11}
\]

从 \(R\) 积分到 \(x\)：

\[
\Delta U(x)
=\int_R^x\mu\omega^2(x'-R)\,dx'
=\frac12\mu\omega^2(x-R)^2. \tag{12}
\]

所以

\[
\boxed{\Delta U=\frac12\mu\omega^2(x-R)^2}. \tag{13}
\]

\needspace{3\baselineskip}
\textbf{（三）解离条件}

Morse 势在平衡点

\[
U(R)=D(1-2)=-D. \tag{14}
\]

当 \(x\to\infty\) 时，

\[
e^{-\alpha(x-R)}\to0,
\qquad U(x)\to0. \tag{15}
\]

因此 \(E=0\) 是刚好能够逃向无穷远的解离阈值。题设
\(E=-D+\varepsilon\)，分子键不解离要求

\[
-D+\varepsilon<0. \tag{16}
\]

振动激发能又不能为负，因此

\[
\boxed{0\le\varepsilon<D}. \tag{17}
\]

\needspace{3\baselineskip}
\textbf{（四）平衡点附近展开}

令

\[
q=x-R. \tag{18}
\]

使用二阶展开

\[
e^{-\alpha q}
=1-\alpha q+\frac12\alpha^2q^2+O(q^3), \tag{19}
\]

\[
e^{-2\alpha q}
=1-2\alpha q+2\alpha^2q^2+O(q^3). \tag{20}
\]

代入 Morse 势：

\[
\begin{aligned}
U
&=D\left[
1-2\alpha q+2\alpha^2q^2
-2\left(1-\alpha q+\frac12\alpha^2q^2\right)
\right]+O(q^3)\\
&=D[-1+\alpha^2q^2]+O(q^3)\\
&=-D+D\alpha^2q^2+O(q^3).
\end{aligned} \tag{21}
\]

与

\[
U\simeq-D+\frac12\mu\omega^2q^2 \tag{22}
\]

比较：

\[
\frac12\mu\omega^2=D\alpha^2. \tag{23}
\]

因此

\[
\boxed{
\omega=\alpha\sqrt{\frac{2D}{\mu}}
}. \tag{24}
\]

振幅为 \(A\) 时，转折点处动能为零，小振动能量高于势阱底的部分为

\[
\varepsilon=D\alpha^2A^2. \tag{25}
\]

所以

\[
\boxed{
A=\frac1\alpha\sqrt{\frac{\varepsilon}{D}}
}. \tag{26}
\]

\needspace{3\baselineskip}
\textbf{（五）一氧化碳的数值计算}

折合质量

\[
\begin{aligned}
\mu
&=\frac{12.00\times15.99}{12.00+15.99}u\\
&=6.8553u\\
&\simeq1.138\times10^{-26}\,\mathrm{kg}.
\end{aligned} \tag{27}
\]

由式（24）

\[
\alpha=\omega\sqrt{\frac{\mu}{2D}}. \tag{28}
\]

代入数值：

\[
\begin{aligned}
\alpha
&=4.088\times10^{14}
\sqrt{
\frac{1.138\times10^{-26}}
{2\times1.780\times10^{-18}}
}\\
&\simeq2.31\times10^{10}\,\mathrm{m^{-1}}.
\end{aligned} \tag{29}
\]

相应的特征长度为

\[
\alpha^{-1}\simeq4.33\times10^{-11}\,\mathrm m. \tag{30}
\]

吸收光子的角频率取 \(\omega\)，其真空波长

\[
\lambda=\frac{2\pi c}{\omega}. \tag{31}
\]

代入 \(c=2.998\times10^8\,\mathrm{m/s}\)：

\[
\lambda
\simeq4.61\times10^{-6}\,\mathrm m
=4.61\,\mu\mathrm m. \tag{32}
\]

故它属于

\[
\boxed{\text{中红外波段}}. \tag{33}
\]

\needspace{3\baselineskip}
\textbf{（六）同位素思考题}

在近似相同的化学势能曲线上，把氕换成质量更大的氘会增大折合质量。由

\[
\omega=\sqrt{\frac{k_{\rm eff}}{\mu}}
\]

可知 \(\mathrm{C-D}\) 键的振动频率低于 \(\mathrm{C-H}\)
键。量子力学中振动基态仍有零点能

\[
E_0=\frac12\hbar\omega.
\]

本题的经典 Morse
模型本身只给振动规律，不能直接推出卤化反应速率。若额外假设反应的控制步骤是断裂
C---H 或 C---D
键、两者电子势能面近似相同，且过渡态的零点能差较小，则基态 C---H
的零点能较高，相应有效能垒较低，预期 C---H
断裂较快。这个条件判断不能推广为所有同位素取代反应都更容易或更困难。

\needspace{3\baselineskip}
\textbf{补充推导与适用条件}

折合质量也可从动力学直接得到。两原子受到大小相等、方向相反的内力
\(F(x)\)：

\[
m_1\ddot x_1=F,
\qquad
m_2\ddot x_2=-F.
\]

相减得

\[
\ddot x=\ddot x_1-\ddot x_2
=F\left(\frac1{m_1}+\frac1{m_2}\right)
=\frac{F}{\mu}. \tag{34}
\]

因此
\(\mu^{-1}=m_1^{-1}+m_2^{-1}\)。这与动能分解得到的式（6）一致，是折合质量的第二条来源。若
\(m_2\gg m_1\)，则 \(\mu\to m_1\)，重原子近似不动；若 \(m_1=m_2=m\)，则
\(\mu=m/2\)，两原子各移动相对距离的一半。

Morse 势可重写为

\[
U(x)=D\left(e^{-\alpha q}-1\right)^2-D,
\qquad q=x-R. \tag{35}
\]

这样势阱底 \(-D\)、无穷远零点和非负平方项一目了然。给定能量
\(E=-D+\varepsilon\)，转折点满足

\[
\left(e^{-\alpha q}-1\right)^2=\frac{\varepsilon}{D}. \tag{36}
\]

令 \(s=\sqrt{\varepsilon/D}\)，在 \(0\le s<1\) 时有两个转折点：

\[
q_{\rm left}=-\frac1\alpha\ln(1+s),
\qquad
q_{\rm right}=-\frac1\alpha\ln(1-s). \tag{37}
\]

左侧转折点为负，右侧为正，而且
\(q_{\rm right}>|q_{\rm left}|\)，准确反映图 5-2
中``压缩侧陡、伸长侧缓''的非对称性。当 \(s\ll1\)
时，\(\ln(1\pm s)\simeq\pm s\)，两侧才近似对称，并恢复

\[
A\simeq\frac{s}{\alpha}
=\frac1\alpha\sqrt{\frac\varepsilon D}. \tag{38}
\]

所以式（26）是小振动振幅，不应在 \(\varepsilon\) 接近 \(D\)
时当成精确转折距离。事实上 \(\varepsilon\to D^-\) 时
\(q_{\rm right}\to\infty\)，而式（26）仍给有限
\(1/\alpha\)，正好显示谐振近似失效。

数值结果可做三重检查。首先 \(\alpha^{-1}\simeq0.433\,\text{Å}\)
小于平衡键长 \(R=1.128\,\text{Å}\)，势能变化尺度合理。其次

\[
\hbar\omega\simeq4.31\times10^{-20}\,\mathrm J
\simeq0.269\,\mathrm{eV},
\]

明显小于键深
\(D\simeq11.1\,\mathrm{eV}\)，低振动态使用谐振近似合理。最后
\(\lambda=2\pi c/\omega\) 的量纲为长度，\(4.61\,\mu\mathrm m\)
位于中红外，不可能是可见光。

\needspace{3\baselineskip}
\textbf{有限能量下的振幅与周期}

取 \(s=\sqrt{\varepsilon/D}\)，\(0<s<1\)。由 \(U=-D+\varepsilon\)
得两个转折点

\[
q_L=-\frac{\ln(1+s)}\alpha,\qquad
q_R=-\frac{\ln(1-s)}\alpha.
\]

因此若把振幅定义为左右转折距离的一半，精确结果是

\[
A_{\rm half}=\frac{1}{2\alpha}\ln\frac{1+s}{1-s},\qquad
q_c=-\frac{1}{2\alpha}\ln(1-s^2).
\]

它们分别是运动区间的半宽和中心偏移。小能量时
\(A_{\rm half}\simeq s/\alpha\)，才回到题目要求的小振动振幅；图5-2分别标出左右距离，避免把不对称振动画成等幅正弦。

有限能量的轨迹还能写为

\[
q(t)=\frac1\alpha\ln\frac{1+s\cos(\Omega t)}{1-s^2},
\qquad
\Omega=\alpha\sqrt{\frac{2(D-\varepsilon)}\mu}.
\]

把 \(e^{-\alpha q}=(1-s^2)/(1+s\cos\Omega t)\) 与 \(\dot q\)
代回能量式可验证。于是 \(T=2\pi/\Omega\)；当 \(\varepsilon\to0\)
回到小振动周期，当 \(\varepsilon\to D^-\)
周期发散，右转折点也趋于无穷。给定 CO 参数满足左转折点仍有
\(x=R+q_L>0\)，没有触及距离坐标的边界。

\subsection{六、关键洞察}\label{ux516dux5173ux952eux6d1eux5bdf-4}

解析图中Morse曲线与谐振抛物线只在谷底接近；有限能量水平线的左右交点不对称。后面的转折点、周期与能量关系据此计算。

\subsection{七、易错点}\label{ux4e03ux6613ux9519ux70b9-4}

\begin{enumerate}
\def\labelenumi{\arabic{enumi}.}
\tightlist
\item
  用 \(m_1+m_2\) 代替折合质量。
\item
  展开指数时漏掉二阶项或把 \(e^{-2\alpha q}\) 的二阶系数写错。
\item
  认为解离阈值是 \(E=D\)；本题势能零点在无穷远，阈值是 \(E=0\)。
\item
  把角频率直接写成 \(c/\lambda\)，漏掉 \(2\pi\)。
\item
  把 \(4.61\,\mu\mathrm m\) 判断为可见光。
\item
  把式（26）当成 Morse 势任意能量下的精确振幅，忽略左右转折点不对称。
\item
  计算吸收光子能量时把角频率代入 \(E=h\omega\)；正确为
  \(E=\hbar\omega\)，等价波长才是 \(2\pi c/\omega\)。
\item
  折合质量用原子质量数直接代入 SI 公式，漏乘原子质量单位 \(u\)。
\item
  把势阱深度 \(D\) 与从势能零点量起的总机械能 \(E\) 混为同一个量。
\end{enumerate}

\subsection{八、课堂讲解建议}\label{ux516bux8bfeux5802ux8bb2ux89e3ux5efaux8bae-4}

在势能图上同时标出 \(-D\)、\(0\) 和
\(-D+\varepsilon\)，让学生先凭图判断解离条件。推导频率时强调``一阶项消失是平衡点，二阶项决定劲度''。

建议先从质心坐标

\[
X=\frac{m_1x_1+m_2x_2}{m_1+m_2}
\]

与相对坐标 \(x=x_1-x_2\)
开始，把总动能完整分成质心平动和相对运动两项。题目位于质心系，所以第一项为零。这样学生会明白折合质量不是经验替换，而是坐标变换后的惯性系数。

第二问可反向教学：先从给定的正弦解两次求导，得到
\(\ddot x=-\omega^2(x-R)\)；再乘 \(\mu\)
识别回复力，最后积分得到势能。每一步都说明从``运动形式''如何反推出``力''和``势能''，训练逆向建模能力。

讲 Morse
势时直接使用式（35）画图，并在同一水平能量线上标出式（37）的两个转折点。先让学生观察右转折点离
\(R\) 更远，再做对数展开。小振动频率只需要二阶系数，解离条件则必须看
\(x\to\infty\)，从而直观区分局部与全局信息。

数值计算建议分行处理：先算 \(\mu/u\) 的无量纲质量数，再乘
\(u\)；再算根号 \(\sqrt{\mu/(2D)}\) 的数量级，最后乘
\(\omega\)。吸收波长可用 \(\nu=\omega/(2\pi)\) 再算
\(c/\nu\)，与直接式（31）互检。

教师验收时应区分三种``能量''：势阱最低值为 \(-D\)，相对势阱底的激发能为
\(\varepsilon\)，从无穷远零点量起的总能量为
\(E=-D+\varepsilon\)。学生若说``解离需要能量大于
\(D\)''，必须追问其零点；正确意思是从势阱最低点至少补充
\(D\)，等价于总能量达到零。统一零点后，式（17）不会产生歧义。

\subsection{九、变式训练}\label{ux4e5dux53d8ux5f0fux8badux7ec3-4}

把 \(^{12}\mathrm C^{16}\mathrm O\) 中的碳替换为
\(^{13}\mathrm C\)，在势能参数近似不变时估算振动频率和吸收波长的相对变化。

势能参数 \(D,\alpha\) 不变时

\[
\omega\propto\mu^{-1/2},
\qquad
\lambda\propto\mu^{1/2}. \tag{39}
\]

原分子的质量数折合量为

\[
\mu_{12}=\frac{12\times16}{12+16}u,
\]

替换后

\[
\mu_{13}=\frac{13\times16}{13+16}u.
\]

因此

\[
\frac{\omega_{13}}{\omega_{12}}
=\sqrt{\frac{\mu_{12}}{\mu_{13}}}<1,
\qquad
\frac{\lambda_{13}}{\lambda_{12}}
=\sqrt{\frac{\mu_{13}}{\mu_{12}}}>1. \tag{40}
\]

重同位素使振动频率降低、吸收波长增大。还可令一个原子质量趋于无穷，验证频率只由较轻原子的质量控制；或要求使用式（37）比较高能振动时左右转折距离，以突出非简谐性。

\newpage

\section{第 6 题}\label{ux7b2c-6-ux9898}

透明介质的相对介电常数和相对磁导率分别为 \(\varepsilon_r,\mu_r\)，折射率
\(n=\sqrt{\varepsilon_r\mu_r}\)。平面电磁波满足

\[
\boldsymbol B=\frac1\omega\boldsymbol k\times\boldsymbol E.
\]

考虑 \(\mu_r=1\) 的两种透明介质，其相对介电常数分别为
\(\varepsilon_{r1},\varepsilon_{r2}\)。一束 \(P\) 偏振光从介质 1
入射到平面界面，入射角、反射角均为 \(i\)，折射角为 \(\gamma\)。

（1）写出界面上关于
\(\boldsymbol E,\boldsymbol D,\boldsymbol H,\boldsymbol B\)
的四个边界连续条件。

（2）推导折射定律。

（3）推导 \(P\) 光的振幅反射比 \(r_p=E_2/E_1\) 和振幅透射比
\(t_p=E_3/E_1\)；说明小角入射且 \(\varepsilon_{r2}>\varepsilon_{r1}\)
时的半波损失，并讨论掠入射。

\begin{figure}[H]
\centering
\includegraphics[width=0.526\linewidth,height=\textheight,keepaspectratio,alt={图6-1　P 偏振光在两介质界面的入射、反射与折射}]{../figures/question-06.pdf}
\caption*{图6-1　P 偏振光在两介质界面的入射、反射与折射}
\end{figure}

\subsection{一、题型定位}\label{ux4e00ux9898ux578bux5b9aux4f4d-5}

先规定三束光的偏振基矢，再投影边界条件。反射系数的符号要和基矢方向一起解释；正入射与掠入射应分别比较实际占主导的场分量。

\subsection{二、知识点}\label{ux4e8cux77e5ux8bc6ux70b9-5}

\begin{enumerate}
\def\labelenumi{\arabic{enumi}.}
\tightlist
\item
  无自由面电荷、无自由面电流时，\(D_n,B_n,E_t,H_t\) 连续。
\item
  界面上任意位置、任意时刻相位相等，要求切向波矢相同。
\item
  对非磁介质，\(k=n\omega/c\)，且 \(B=nE/c\)。
\item
  \(P\) 光的电场在入射面内，磁场垂直入射面。
\item
  反射系数的正负必须与所选偏振基矢约定一起解释。
\end{enumerate}

\subsection{三、物理图像}\label{ux4e09ux7269ux7406ux56feux50cf-5}

界面不能凭空积累无限大的电荷或电流，因此电磁场的法向、切向分量受到连续性约束。反射波的传播方向改变后，为保持
\(\boldsymbol E,\boldsymbol B,\boldsymbol k\)
的右手关系，其偏振基矢方向也随之改变；``系数正负''和``实验室固定方向上的相位翻转''不能混为一谈。

\subsection{四、解题入口}\label{ux56dbux89e3ux9898ux5165ux53e3-5}

先写矢量边界条件，再由界面相位匹配得到 Snell
定律。振幅系数推导中严格按图定义三个电场振幅为沿各自箭头方向的标量，投影到界面切向和法向后联立。

\subsection{五、解析}\label{ux4e94ux89e3ux6790-5}

\begin{figure}[H]
\centering
\includegraphics[width=0.897\linewidth,height=\textheight,keepaspectratio,alt={图6-2　P 光电场分量、偏振基矢和反射相位约定}]{../figures/solution-06.pdf}
\caption*{图6-2　P 光电场分量、偏振基矢和反射相位约定}
\end{figure}

图6-2沿每束光分别规定正电场基矢，反射系数曲线按此约定穿过布儒斯特零点并趋于
\(-1\)。零振幅处的相位和固定坐标分量的相位要另作判断。

\needspace{3\baselineskip}
\textbf{（一）四个边界条件}

设单位法向 \(\hat{\boldsymbol n}\) 从介质 1 指向介质
2。若界面没有自由面电荷 \(\sigma_f\) 和自由面电流
\(\boldsymbol K_f\)，则

\[
\boxed{
\hat{\boldsymbol n}\cdot
(\boldsymbol D_2-\boldsymbol D_1)=0
}. \tag{1}
\]

\[
\boxed{
\hat{\boldsymbol n}\cdot
(\boldsymbol B_2-\boldsymbol B_1)=0
}. \tag{2}
\]

\[
\boxed{
\hat{\boldsymbol n}\times
(\boldsymbol E_2-\boldsymbol E_1)=\boldsymbol0
}. \tag{3}
\]

\[
\boxed{
\hat{\boldsymbol n}\times
(\boldsymbol H_2-\boldsymbol H_1)=\boldsymbol0
}. \tag{4}
\]

实际应用时，介质 1 一侧的总场是``入射场加反射场''，介质 2 一侧是透射场。

\needspace{3\baselineskip}
\textbf{（二）相位匹配与折射定律}

在界面上选切向坐标 \(x\)。三列波的相位分别包含

\[
k_1x\sin i,\qquad
k_1x\sin i,\qquad
k_2x\sin\gamma. \tag{5}
\]

为了使边界条件在任意 \(x,t\)
都成立，三列波必须有相同频率，且切向波矢相同：

\[
k_1\sin i=k_2\sin\gamma. \tag{6}
\]

透明介质中

\[
k_j=\frac{n_j\omega}{c}. \tag{7}
\]

代入并约去 \(\omega/c\)：

\[
\boxed{
n_1\sin i=n_2\sin\gamma
}. \tag{8}
\]

由于 \(\mu_{r1}=\mu_{r2}=1\)，

\[
n_j=\sqrt{\varepsilon_{rj}}. \tag{9}
\]

故也可写成

\[
\sqrt{\varepsilon_{r1}}\sin i
=\sqrt{\varepsilon_{r2}}\sin\gamma. \tag{10}
\]

\needspace{3\baselineskip}
\textbf{（三）P 光的分量方程}

按题图箭头方向定义入射、反射、透射电场振幅分别为
\(E_1,E_2,E_3\)。切向电场连续给出

\[
E_1\cos i-E_2\cos i=E_3\cos\gamma. \tag{11}
\]

法向电位移连续给出

\[
\varepsilon_{r1}E_1\sin i
+\varepsilon_{r1}E_2\sin i
=\varepsilon_{r2}E_3\sin\gamma. \tag{12}
\]

由 Snell 定律

\[
\frac{\varepsilon_{r2}}{\varepsilon_{r1}}
=\frac{n_2^2}{n_1^2}
=\frac{\sin^2 i}{\sin^2\gamma}. \tag{13}
\]

由式（11）先把透射振幅写成

\[
E_3=(E_1-E_2)\frac{\cos i}{\cos\gamma}. \tag{14}
\]

令

\[
r_p=\frac{E_2}{E_1}.
\]

则式（14）除以 \(E_1\) 后为

\[
\frac{E_3}{E_1}=(1-r_p)\frac{\cos i}{\cos\gamma}. \tag{15}
\]

把它代入式（12），并约去 \(E_1\)：

\[
\varepsilon_{r1}(1+r_p)\sin i
=\varepsilon_{r2}(1-r_p)
\frac{\cos i}{\cos\gamma}\sin\gamma. \tag{16}
\]

右端的几何因子整理为 \(\cos i\tan\gamma\)，再把式（13）代入：

\[
\begin{aligned}
1+r_p
&=\frac{\varepsilon_{r2}}{\varepsilon_{r1}}
\frac{\cos i\tan\gamma}{\sin i}(1-r_p)\\
&=\frac{\sin^2i}{\sin^2\gamma}
\frac{\cos i\sin\gamma}
{\cos\gamma\sin i}(1-r_p)\\
&=\frac{\sin i\cos i}
{\sin\gamma\cos\gamma}(1-r_p).
\end{aligned} \tag{17}
\]

两端乘以 \(\sin\gamma\cos\gamma\)，把含 \(r_p\) 的项移到一侧：

\[
r_p\bigl(\sin i\cos i+\sin\gamma\cos\gamma\bigr)
=\sin i\cos i-\sin\gamma\cos\gamma. \tag{18}
\]

因此

\[
r_p=
\frac{\sin i\cos i-\sin\gamma\cos\gamma}
{\sin i\cos i+\sin\gamma\cos\gamma}. \tag{19}
\]

利用积化和差恒等式

\[
\sin i\cos i-\sin\gamma\cos\gamma
=\sin(i-\gamma)\cos(i+\gamma),
\]

\[
\sin i\cos i+\sin\gamma\cos\gamma
=\sin(i+\gamma)\cos(i-\gamma),
\]

得到

\[
\boxed{
r_p\equiv\frac{E_2}{E_1}
=\frac{\tan(i-\gamma)}{\tan(i+\gamma)}
}. \tag{20}
\]

再由式（11）

\[
t_p\equiv\frac{E_3}{E_1}
=\frac{(1-r_p)\cos i}{\cos\gamma}. \tag{21}
\]

由式（19）可得

\[
1-r_p
=\frac{2\sin\gamma\cos\gamma}
{\sin(i+\gamma)\cos(i-\gamma)}. \tag{22}
\]

代入上式并约去 \(\cos\gamma\)：

\[
\boxed{
t_p=
\frac{2\sin\gamma\cos i}
{\sin(i+\gamma)\cos(i-\gamma)}
}. \tag{23}
\]

也可写成折射率形式

\[
t_p=
\frac{2n_1\cos i}
{n_2\cos i+n_1\cos\gamma}, \tag{24}
\]

但要注意式（24）所用的偏振基矢符号约定。

\needspace{3\baselineskip}
\textbf{（四）小角入射的半波损失}

当 \(i,\gamma\ll1\) 时，

\[
\tan(i\pm\gamma)\simeq i\pm\gamma. \tag{25}
\]

若 \(\varepsilon_{r2}>\varepsilon_{r1}\)，则 \(n_2>n_1\)。由 Snell
定律，小角下

\[
n_1i\simeq n_2\gamma,
\]

所以

\[
i>\gamma. \tag{26}
\]

于是式（20）的标量
\(r_p>0\)。但按题图对反射波所选的偏振基矢，反射电场箭头与入射电场在近法线方向上相反。因此实验室固定横向方向的电场发生符号翻转，即相位增加
\(\pi\)，这就是光从光疏介质射向光密介质时的半波损失。

在严格正入射极限，用固定实验室横向方向定义反射系数更直观：

\[
r_{\rm lab}
=\frac{n_1-n_2}{n_1+n_2}<0. \tag{27}
\]

\needspace{3\baselineskip}
\textbf{（五）布儒斯特角和掠入射}

当

\[
i+\gamma=\frac{\pi}{2}, \tag{28}
\]

式（20）的分母发散而 \(r_p=0\)，这就是 \(P\) 光的布儒斯特角。结合 Snell
定律可得

\[
\tan i_B=\frac{n_2}{n_1}. \tag{29}
\]

这里先取 \(n_2>n_1\)，透射角仍为实数。当 \(i\to\pi/2\)
时，\(r_p\to-1\)，反射率趋近 \(1\)。按图中基矢，切向电场比为
\(E_{rx}/E_{ix}=-r_p\to+1\)，法向电场比为
\(E_{rz}/E_{iz}=r_p\to-1\)。掠入射时电场主要沿法向，入射与反射偏振基矢趋于同向，主分量发生相位反转；不能给所有实验室分量统一贴上同一个符号。

若
\(n_1>n_2\)，在达到掠入射之前已经进入全反射区，透射角不能继续作为实角代入上述三角式；应使用复数传播方向，反射振幅模为1、相位随角度变化。布儒斯特角处反射场为零，它的相位没有定义。

\needspace{3\baselineskip}
\textbf{补充推导与适用条件}

四个边界条件各有明确积分来源。用跨越界面的高斯小柱，令高度趋于零，可由
\(\nabla\cdot\boldsymbol D=\rho_f\) 得

\[
\hat{\boldsymbol n}\cdot(\boldsymbol D_2-\boldsymbol D_1)=\sigma_f;
\]

无自由面电荷时右端为零。由 \(\nabla\cdot\boldsymbol B=0\) 同理得到
\(B_n\) 连续。再取跨界面的极窄矩形回路，法拉第定律在面积趋零时给出
\(E_t\) 连续；安培---麦克斯韦定律给出

\[
\hat{\boldsymbol n}\times(\boldsymbol H_2-\boldsymbol H_1)=\boldsymbol K_f.
\]

无自由面电流时 \(H_t\)
连续。这里连续的是指定分量，不是四个矢量整体逐分量都连续。

折射定律来自界面上任一点、任一时刻的相位必须匹配。若切向波矢不同，边界条件中的三列指数函数会随
\(x\) 以不同周期变化，不可能对所有 \(x\) 同时成立。因此

\[
\omega_1=\omega_2=\omega_3,
\qquad
\boldsymbol k_{1\parallel}
=\boldsymbol k_{2\parallel}
=\boldsymbol k_{3\parallel}. \tag{30}
\]

反射角等于入射角和 Snell
定律都是这一矢量相位条件的结果，而不是几何光学的独立经验公设。

能流提供重要检查。非磁介质中，平面波平均坡印廷矢量大小正比于
\(nE^2\)。若按同一振幅约定，能流反射率和透射率应写成

\[
\mathcal R=|r_p|^2,
\qquad
\mathcal T=\frac{n_2\cos\gamma}{n_1\cos i}|t_p|^2. \tag{31}
\]

无吸收界面必须满足

\[
\mathcal R+\mathcal T=1. \tag{32}
\]

在布儒斯特角 \(r_p=0\)，于是 \(\mathcal T=1\)；掠入射时
\(\mathcal R\to1\)、\(\mathcal T\to0\)。这两个极限同时检查振幅公式和能流几何因子。

正入射处 \(P\)、\(S\) 偏振应给相同物理反射。若反射波的 \(P\)
偏振基矢为了保持 \((\boldsymbol E,\boldsymbol B,\boldsymbol k)\)
右手关系而反向，标量 \(r_p\)
的符号会与固定实验室横向分量的符号相反。教师应把``基矢变化''画在图 6-2
上，而不是只用文字解释半波损失。

式（20）在 \(i+\gamma=\pi/2\)
时应理解为极限：分子有限、分母的正切发散，故
\(r_p\to0\)。不能把``分母无穷大''当作公式无意义。掠入射同样应使用极限而不是直接把
\(\cos i=0\) 代入中间分量方程后消去信息。

\needspace{3\baselineskip}
\textbf{不用正切奇点的系数与图上投影}

同一基矢约定下，把边界方程写成折射率形式更便于计算：

\[
r_p=\frac{n_2\cos i-n_1\cos\gamma}{n_2\cos i+n_1\cos\gamma},
\qquad
 t_p=\frac{2n_1\cos i}{n_2\cos i+n_1\cos\gamma}.
\]

这两式在正入射处可直接取极限，避免正切形式的 \(0/0\)。图6-1中的波矢为
\((\sin i,-\cos i)\)、\((\sin i,\cos i)\)、\((\sin\gamma,-\cos\gamma)\)，相应电场基矢为
\((\cos i,\sin i)\)、\((-\cos i,\sin i)\)、\((\cos\gamma,\sin\gamma)\)，每一对都垂直。

图6-2按 \(n_2/n_1=1.5\) 绘出 \(r_p\) 和
\(\mathcal R\)，不是任意示意曲线：\(r_p(0)=0.2\)，到 \(i_B=\arctan1.5\)
为零，随后降至 \(-1\)；\(\mathcal R\)
则从0.04降到0，再升到1。图中的系数是沿各自偏振基矢的标量，判断实验室分量相位仍须使用上面的投影。

\needspace{3\baselineskip}
\textbf{反射系数的符号与可测相位}

让学生先在三条光线上画出与波矢垂直的正电场基矢，再写
\(r_p\)。若把反射基矢整体翻转，标量 \(r_p\)
也随之变号，但实验室中的电场和反射率不变。给相位结论时必须同时报出所用基矢，不能从一个孤立的负号推断所有分量都反相。

在布儒斯特角处反射振幅为零，零复数没有唯一相位。可以比较零点两侧非零振幅的相位跳变，却不能说零点处相位恰为某个角。掠入射极限则应分别投影切向与法向电场，再检查法向能流；全反射区需要复数折射角，实角正切式不可直接沿用。

图6-2的 \(r_p\) 与 \(\mathcal R=r_p^2\)
来自同一组折射率。课堂追问可以要求指出蓝线的负值为何没有让红线变负；答案应涉及振幅的符号、强度的平方和能流系数三个不同层次。

\subsection{六、关键洞察}\label{ux516dux5173ux952eux6d1eux5bdf-5}

解析图沿每束光分别规定正电场基矢，反射系数曲线按此约定穿过布儒斯特零点并趋于
\(-1\)。零振幅处的相位和固定坐标分量的相位要另作判断。

\subsection{七、易错点}\label{ux4e03ux6613ux9519ux70b9-5}

\begin{enumerate}
\def\labelenumi{\arabic{enumi}.}
\tightlist
\item
  把边界条件写成 \(\boldsymbol E,\boldsymbol B\) 四个矢量全部连续。
\item
  忘记介质 1 一侧是入射场与反射场之和。
\item
  折射定律把 \(k\) 与真空波长混用。
\item
  投影时漏掉反射电场切向分量的负号。
\item
  不说明符号约定便直接用``\(r_p\) 正或负''判断相位。
\item
  法向 \(D\) 连续条件中只写透射场，漏掉介质 1
  侧入射与反射的法向分量代数和。
\item
  把振幅透射率 \(|t_p|^2\) 直接当能量透射率，漏掉 \(n\cos\theta\)
  的能流因子。
\item
  在布儒斯特角把 \(\tan(i+\gamma)\) 的发散误判成反射系数发散。
\item
  正入射时仍坚持使用没有定义好的入射面，从而人为制造 \(P\)、\(S\)
  结果不一致。
\end{enumerate}

\subsection{八、课堂讲解建议}\label{ux516bux8bfeux5802ux8bb2ux89e3ux5efaux8bae-5}

先固定坐标和三个偏振基矢，再写每个场的切向、法向分量表格。推导完成后分别检查正入射、布儒斯特角和掠入射三个极限，这比只记公式更可靠。

板书第一行必须画界面法向、三条波矢和三个 \(P\)
偏振电场箭头，并规定每个标量振幅沿哪一支基矢为正。第二行做分量表：\(E_t,E_n,H_t,B_n\)。只有符号表完成后才写边界方程，可避免在式（11）、式（12）中凭直觉加减。

折射定律建议从相位写起。在界面上取两个任意点，比较同一时刻相位差，说明切向波矢必须相等；这比直接引用
Snell 定律更符合题目``推导''的要求。接着用
\(\boldsymbol B=(1/\omega)\boldsymbol k\times\boldsymbol E\)
确定磁场方向和幅值，再选择两条独立边界条件求 \(r_p,t_p\)。

得到公式后按三个极限组织复核。正入射要换到固定实验室方向得到
\((n_1-n_2)/(n_1+n_2)\)；布儒斯特角检查反射为零；掠入射检查几乎全反射。最后代式（31）验证能流守恒，使学生看到振幅与能量系数的区别。

课堂验收不接受只背 Fresnel
公式而不写符号约定。完整答案至少应说明：介质编号、法向、偏振基矢、反射振幅定义，以及``半波损失''比较的是同一实验室方向的电场。

建议让学生把式（20）与折射率形式互相化简一次。化简过程中只使用 Snell
定律和和差公式；完成后在正入射、布儒斯特角各代一次。两种形式数值相同却可能因基矢约定呈现不同标量符号，这正适合用来强调``公式外观''不等于``物理相位约定''。最后用式（32）验证能流，形成相位---振幅---能量三个层次的闭环。

\subsection{九、变式训练}\label{ux4e5dux53d8ux5f0fux8badux7ec3-5}

用同样方法推导 \(S\) 偏振光的振幅反射系数，并说明为什么 \(S\)
光不存在布儒斯特角。

\(S\) 光电场垂直入射面，切向 \(E\) 连续给

\[
E_i+E_r=E_t.
\]

切向 \(H\) 连续并利用非磁介质波阻抗，可得

\[
r_s=\frac{n_1\cos i-n_2\cos\gamma}
{n_1\cos i+n_2\cos\gamma}. \tag{33}
\]

若 \(n_1,n_2>0\) 且两角都在传播区间，分母为正。要使分子为零需
\(n_1\cos i=n_2\cos\gamma\)，再与 \(n_1\sin i=n_2\sin\gamma\)
联立，只能得到 \(n_1=n_2\)
的平凡无界面情形。因此普通两介质界面不存在非平凡的 \(S\)
光布儒斯特角。可进一步验证正入射时
\(r_s=(n_1-n_2)/(n_1+n_2)\)，并用能流系数检查守恒。

\newpage

\section{第 7 题}\label{ux7b2c-7-ux9898}

实验室有一节电池，电动势 \(\varepsilon\) 约为 \(1.5\,\mathrm V\)、内阻
\(r\) 约为 \(1.0\,\Omega\)；一只电流表量程 \(I_g=20\,\mathrm{mA}\)、内阻
\(r_g\) 约为 \(20\,\Omega\)；另有定值电阻 \(R_0=20\,\Omega\) 和待测电阻
\(R_x\) 约为
\(100\,\Omega\)。提供两个单刀双掷开关及导线。要求电路接好后不再改线，只通过切换开关测出
\(\varepsilon,r,R_x,r_g\)。

（1）画出电路图。

（2）简述测量步骤。

（3）用各次读数及已知 \(R_0\) 表示四个未知量；允许用二次方程表示。

\begin{figure}[H]
\centering
\includegraphics[width=0.763\linewidth,height=\textheight,keepaspectratio,alt={图7-1　用两个单刀双掷开关实现四种固定接线状态}]{../figures/question-07.pdf}
\caption*{图7-1　用两个单刀双掷开关实现四种固定接线状态}
\end{figure}

\subsection{一、题型定位}\label{ux4e00ux9898ux578bux5b9aux4f4d-6}

先把两只开关的公共端和触点接清楚，再追踪四种电流路径。参考方案需要第一只刀开关断开；两种串联读数确定电动势与总电阻，分流读数再拆出表头和电池内阻。

\subsection{二、知识点}\label{ux4e8cux77e5ux8bc6ux70b9-6}

\begin{enumerate}
\def\labelenumi{\arabic{enumi}.}
\tightlist
\item
  电池内阻和电流表内阻都必须计入闭合电路。
\item
  参考方案用到第一只刀开关断开的位置，两只开关共同形成四种测量状态。
\item
  先用两组纯串联读数求电动势与总未知电阻和。
\item
  再用两种并联分流状态分别确定 \(r_g\) 和 \(r\)。
\item
  二次方程的两个数学根要结合估计值和量程选择物理解。
\end{enumerate}

\subsection{三、物理图像}\label{ux4e09ux7269ux7406ux56feux50cf-6}

前两种状态只改变是否串入
\(R_0\)，因此同一电源电动势可以被直接消去。第三种状态让电流表与 \(R_0\)
并联，分流比暴露出 \(r_g\)；第四种状态让``\(R_x+r_g\) 支路''与 \(R_0\)
并联，电源总电流与表头支路电流之差暴露出内阻 \(r\)。

\subsection{四、解题入口}\label{ux56dbux89e3ux9898ux5165ux53e3-6}

按图连接后依次记录四个表头读数 \(I_1,I_2,I_3,I_4\)。先由前两式求
\(\varepsilon\) 和

\[
S=R_x+r_g+r.
\]

随后把第三、第四状态分别整理成关于 \(r_g\) 和 \(r\) 的二次方程。

\subsection{五、解析}\label{ux4e94ux89e3ux6790-6}

\begin{figure}[H]
\centering
\includegraphics[width=0.883\linewidth,height=\textheight,keepaspectratio,alt={图7-2　四种开关状态的等效电路及电流分配}]{../figures/solution-07.pdf}
\caption*{图7-2　四种开关状态的等效电路及电流分配}
\end{figure}

图7-2把四个开关状态化成四个闭合回路；图7-1标出a、b、c、d结点和两公共端。四个读数按同一触点表列式，根的筛选再利用题给电阻范围。

\needspace{3\baselineskip}
\textbf{（一）四种测量状态}

按图中编号切换开关，得到：

\begin{enumerate}
\def\labelenumi{\arabic{enumi}.}
\tightlist
\item
  状态 A：\(R_x\)、电流表、\(R_0\) 与电池串联，读数为 \(I_1\)。
\item
  状态 B：\(R_0\) 被旁路，\(R_x\)、电流表与电池串联，读数为 \(I_2\)。
\item
  状态 C：电流表与 \(R_0\) 并联，再与 \(R_x\) 及电池内阻串联，表头读数为
  \(I_3\)。
\item
  状态 D：\(R_x\) 与电流表串联成一支路，该支路与 \(R_0\)
  并联接到电池，表头读数为 \(I_4\)。
\end{enumerate}

每次切换后待读数稳定再记录。应先从总电阻较大的状态 A
开始，以免表头过载。

\needspace{3\baselineskip}
\textbf{（二）状态 A、B：求 \(\varepsilon\)}

状态 A 的闭合电路方程：

\[
\varepsilon
=I_1(R_x+R_0+r_g+r). \tag{1}
\]

状态 B：

\[
\varepsilon
=I_2(R_x+r_g+r). \tag{2}
\]

定义

\[
S=R_x+r_g+r. \tag{3}
\]

则

\[
\varepsilon=I_1(S+R_0)=I_2S. \tag{4}
\]

由

\[
I_1S+I_1R_0=I_2S
\]

得

\[
\boxed{
S=\frac{I_1R_0}{I_2-I_1}
}. \tag{5}
\]

再由 \(\varepsilon=I_2S\)：

\[
\boxed{
\varepsilon
=\frac{R_0I_1I_2}{I_2-I_1}
}. \tag{6}
\]

由于撤去 \(R_0\) 后电流增大，应有 \(I_2>I_1\)。

\needspace{3\baselineskip}
\textbf{（三）状态 C：求电流表内阻 \(r_g\)}

状态 C 中，电流表与 \(R_0\) 两端电压相等。通过 \(R_0\) 的电流为

\[
I_{R_0}=\frac{I_3r_g}{R_0}. \tag{7}
\]

流经 \(R_x\) 和电池内阻的总电流为

\[
I_{\rm tot}
=I_3+I_{R_0}
=I_3\left(1+\frac{r_g}{R_0}\right). \tag{8}
\]

回路电压关系：

\[
\varepsilon
=I_{\rm tot}(R_x+r)+I_3r_g. \tag{9}
\]

由 \(S=R_x+r_g+r\)，

\[
R_x+r=S-r_g. \tag{10}
\]

代入式（8）---式（10）：

\[
\frac{\varepsilon}{I_3}
=\left(1+\frac{r_g}{R_0}\right)(S-r_g)+r_g. \tag{11}
\]

展开：

\[
\begin{aligned}
\frac{\varepsilon}{I_3}
&=S-r_g+\frac{Sr_g-r_g^2}{R_0}+r_g\\
&=S+\frac{r_g(S-r_g)}{R_0}.
\end{aligned} \tag{12}
\]

所以

\[
r_g(S-r_g)
=R_0\left(\frac{\varepsilon}{I_3}-S\right). \tag{13}
\]

整理为二次方程：

\[
\boxed{
r_g^2-Sr_g
+R_0\left(\frac{\varepsilon}{I_3}-S\right)=0
}. \tag{14}
\]

解为

\[
\boxed{
r_g=
\frac{
S\pm\sqrt{
S^2-4R_0\left(\dfrac{\varepsilon}{I_3}-S\right)
}
}{2}
}. \tag{15}
\]

应选择接近题给估计值 \(20\,\Omega\)、并能使全部电阻为正的根。

\needspace{3\baselineskip}
\textbf{（四）状态 D：求电池内阻 \(r\)}

状态 D 中，含表头的支路电阻为

\[
X=R_x+r_g. \tag{16}
\]

该支路电流为 \(I_4\)，故并联端电压为

\[
U=I_4X. \tag{17}
\]

通过 \(R_0\) 的电流为

\[
I_{R_0}=\frac{I_4X}{R_0}. \tag{18}
\]

电池总电流

\[
I_{\rm tot}
=I_4\left(1+\frac{X}{R_0}\right). \tag{19}
\]

电动势等于端电压加内阻压降：

\[
\varepsilon
=I_4X
+I_4\left(1+\frac{X}{R_0}\right)r. \tag{20}
\]

而

\[
S=X+r,\qquad X=S-r. \tag{21}
\]

代入并化简：

\[
\begin{aligned}
\frac{\varepsilon}{I_4}
&=S-r+r+\frac{(S-r)r}{R_0}\\
&=S+\frac{r(S-r)}{R_0}.
\end{aligned} \tag{22}
\]

故

\[
\boxed{
r^2-Sr
+R_0\left(\frac{\varepsilon}{I_4}-S\right)=0
}. \tag{23}
\]

所以

\[
\boxed{
r=
\frac{
S\pm\sqrt{
S^2-4R_0\left(\dfrac{\varepsilon}{I_4}-S\right)
}
}{2}
}. \tag{24}
\]

选取接近 \(1\,\Omega\) 的物理解。最后

\[
\boxed{
R_x=S-r_g-r
}. \tag{25}
\]

至此 \(\varepsilon,r,R_x,r_g\) 均由 \(I_1,I_2,I_3,I_4,R_0\) 表出。

\needspace{3\baselineskip}
\textbf{补充推导与适用条件}

状态 A、B 的作用是先把四个未知量压缩。两式相除也可写成

\[
\frac{I_2}{I_1}=\frac{S+R_0}{S}
=1+\frac{R_0}{S}. \tag{26}
\]

因此

\[
S=\frac{R_0}{I_2/I_1-1},
\]

与式（5）完全一致。这个比值形式显示：若 \(I_2\) 与 \(I_1\)
太接近，分母很小，\(S\) 的测量误差会被显著放大；实验设计中 \(R_0\) 应与
\(S\) 同数量级，而不能远小于所有未知电阻。

状态 C 的二次方程本质上来自

\[
r_g(S-r_g)=K_3,
\qquad
K_3=R_0\left(\frac\varepsilon{I_3}-S\right). \tag{27}
\]

若 \(r_g\) 是一个根，则 \(S-r_g=R_x+r\)
是另一个根。这是二次方程两根之和为 \(S\)
的物理含义。题给估计值告诉我们应选择接近 \(20\,\Omega\) 的小根，而不是约
\(R_x+r\simeq101\,\Omega\) 的大根。状态 D 同理：

\[
r(S-r)=K_4,
\]

两根分别对应 \(r\) 与 \(R_x+r_g\)，应选接近 \(1\,\Omega\)
的小根。根的二义性不是代数错误，而是只凭单个并联状态无法辨认``某一支路电阻''和``其余串联电阻''；先验量级或更多测量负责消除二义性。

用题给估计值预估电流，可以检查电路是否安全。取
\(\varepsilon=1.5\,\mathrm V\)、\(r=1\,\Omega\)、\(r_g=20\,\Omega\)、\(R_x=100\,\Omega\)：

\[
I_1\simeq\frac{1.5}{141}=10.6\,\mathrm{mA},
\qquad
I_2\simeq\frac{1.5}{121}=12.4\,\mathrm{mA}. \tag{28}
\]

状态 C 中
\(r_g\parallel R_0=10\,\Omega\)，表头分得并联总电流的一半，约为
\(6.8\,\mathrm{mA}\)。状态 D 中
\((R_x+r_g)\parallel R_0\simeq17.1\,\Omega\)，含表头支路电流约为
\(11.8\,\mathrm{mA}\)。四个读数均低于
\(20\,\mathrm{mA}\)，也彼此有可分辨差别。

四状态电路必须``接好后不再改线''，所以图 7-1
的两个开关不是四张彼此独立的电路图，而是同一网络的四个拓扑。教师核对重绘图时，应沿每个开关位置实际追踪电流路径，确认不存在短路电池、悬空表头或把
\(R_0\) 接到错误结点的情况。

由实测数据反演后还应做残差检查。把求得的 \(\varepsilon,r,r_g,R_x\)
重新代回四个原始电路，计算预测电流
\(I_i^{\rm pred}\)，再与读数比较。若二次方程选错根，往往仍能满足对应的一条代数方程，却会与估计量级、其他状态或量程约束冲突。

\needspace{3\baselineskip}
\textbf{固定接线、开关位置与禁用状态}

图7-1定义结点：电池正端为 \(a\)，\(R_x\) 后为 \(b\)，表头后为
\(c\)，\(R_0\) 后为 \(d\)，电池负端接 \(K_2\) 的公共端。\(K_1\)
的公共端接 \(d\)，触点1接 \(a\)、触点2接 \(b\)；\(K_2\) 的触点3接
\(d\)、触点4接 \(c\)。

四次操作为：A：\(K_1\) 断开、\(K_2\) 接3；B：\(K_1\) 断开、\(K_2\)
接4；C：\(K_1\) 接2、\(K_2\) 接4；D：\(K_1\) 接1、\(K_2\) 接4。A沿
\(a\to R_x\to b\to A\to c\to R_0\to d\) 返回电池；B的 \(R_0\)
一端悬空；C中 \(d=b\)，故 \(R_0\) 与表头同接 \(b,c\)；D中 \(d=a\)，故
\(R_0\) 与 \(R_x+r_g\) 同接
\(a,c\)。这四条路径对应图7-2及式（1）、（2）、（9）、（20）。

原参考图所写的``空''是第一只刀开关的断开状态，并非只在两个掷端之间切换。采用能够抬刀断开的实验开关时可以实现；若器材是没有断开位置的严格两位置转换开关，这个参考方案不能照搬。状态
\(K_1\) 接1、\(K_2\) 接3会使电池短路，必须排除。测量时先断开
\(K_1\)，再按表切换，避免经过这一组合。

二次方程还可用较稳定的数值形式求小根。若 \(x^2-Sx+K=0\)，则

\[
x_{\rm small}=\frac{2K}{S+\sqrt{S^2-4K}},
\]

避免直接相减两个接近数。\(K_3,K_4\) 必须非负且不超过
\(S^2/4\)，否则读数与无源电阻模型不相容。以题给估计值生成四个未舍入电流再反演，会得到
\(\varepsilon=1.5\,\mathrm V\)、\(r=1\,\Omega\)、\(r_g=20\,\Omega\)、\(R_x=100\,\Omega\)；量程和根筛选应在回代中共同核对。

\needspace{3\baselineskip}
\textbf{逐状态核对电池端口}

讲接线时先把a、b、c、d四结点编号，电池正极恒接a，负极是 \(K_2\)
公共端。逐行读触点表，先检查有没有电池短路，再判断 \(R_0\)
是串联、旁路还是并联。\(K_1\)
的``空''状态表示不接1、2，不能被当作接在悬空导线上但仍改变结点。

可以用 \(R_x=100,r_g=20,r=1,R_0=150\) 欧姆和 \(\varepsilon=1.5\)
伏生成四个理论读数。对并联的状态，电流表读的是自己的支路电流，电池总电流不能直接代入；电池内阻压降仍由总电流决定。先从
\(I_1,I_2\) 求
\(\varepsilon,S\)，再用另外两状态求两二次方程，最后同时筛选
\(r_g\approx20\) 与 \(r\approx1\) 的根。

仪表误差增大时，判别式可能变负。此时应记录读数与模型不相容，检查接线和量程；不应把负判别式直接取绝对值。两个接近根相减时，用稳定公式
\(2K/[S+\sqrt{S^2-4K}]\) 保留小根的精度。

\needspace{3\baselineskip}
\textbf{用未舍入的四个读数核对根筛选}

取本题估计值
\(\varepsilon=1.5\)伏、\(R_x=100\)欧、\(r_g=20\)欧、\(r=1\)欧，并以
\(R_0=100\)欧作演示。四种负载路径直接给

\[
I_1=\frac{1.5}{221},\qquad I_2=\frac{1.5}{121},
\]

\[
I_3=\frac{1.5\,[100/(100+20)]}{101+100\times20/(100+20)},
\]

\[
I_4=\frac{1.5\,[100/(100+100+20)]}{1+100(100+20)/(100+100+20)}.
\]

I3、I4是流过表头的电流，不是电池总电流；分流系数必须在电池端等效阻抗之后再乘。把这些未舍入值回代，得到S=121欧、表头二次式的两根20与101欧、内阻二次式的两根1与120欧。较大根与题给估计范围冲突；选择小根后
\(R_x=S-r_g-r=100\)欧。根筛选来自先验范围，不是二次方程自动只剩一个根。

若去掉估计范围，只给四个电流，就应枚举两组根并检查全部电阻非负、四条读数同时回代以及开关能否实现。某个二次式出现两个正根，并不证明两个完整电路都可行；也不能只因习惯上内阻小便不说明地删根。

从 \(S=I_1R_0/(I_2-I_1)\)
还能看出，当两串联电流很接近时，读数误差会被相减后的分母放大。实验方案宜让R0足够产生可分辨的电流变化，同时保持表头在量程内；R0无限大又会使I1太小，不能只追求读数差大。这个误差讨论不改变原题精确电流模型，实际反算范围还需要各表的误差资料。

\subsection{六、关键洞察}\label{ux516dux5173ux952eux6d1eux5bdf-6}

解析图把四个开关状态化成四个闭合回路；图7-1标出a、b、c、d结点和两公共端。四个读数按同一触点表列式，根的筛选再利用题给电阻范围。

\subsection{七、易错点}\label{ux4e03ux6613ux9519ux70b9-6}

\begin{enumerate}
\def\labelenumi{\arabic{enumi}.}
\tightlist
\item
  状态 A、B 漏计表头内阻或电池内阻。
\item
  并联状态把表头读数误当成电池总电流。
\item
  由电压相等写分流关系时把 \(r_g/R_0\) 写反。
\item
  二次方程两个根不做物理筛选。
\item
  未检查电流是否超过 \(20\,\mathrm{mA}\) 量程。
\item
  把状态 C、D 的二次方程大根也当作可接受解，没有识别两根互为 \(x\) 与
  \(S-x\)。
\item
  用 \(I_3\)、\(I_4\) 作为电池总电流，漏掉与表头并联支路中的电流。
\item
  先对读数过早四舍五入，再用接近的 \(I_1,I_2\) 作差，造成 \(S\) 和
  \(\varepsilon\) 的明显误差。
\item
  只画四张等效图，没有给出一个能由两个单刀双掷开关实际切换得到的统一接线图。
\end{enumerate}

\subsection{八、课堂讲解建议}\label{ux516bux8bfeux5802ux8bb2ux89e3ux5efaux8bae-6}

先不推公式，让学生把四个等效电路分别重画，并在每条支路标电流。确认表头读数在四种状态下分别代表什么，再进入代数消元。

课堂开始先给学生两个任务：任何状态都不得让电池直接短路；任何状态表头都不得超量程。用估计值做式（28）的预估后，再讨论开关拓扑。这样设计不只是代数可解，还兼顾实验安全与信噪比。

推导按``先总和、后拆分''的顺序。状态 A、B 只出现 \(S=R_x+r_g+r\)，先求
\(S,\varepsilon\)；状态 C 只把 \(r_g\) 从 \(S\) 中拆出；状态 D 只把
\(r\) 从 \(S\) 中拆出；最后用差求
\(R_x\)。每完成一个阶段，都把已经确定与仍未知的量列在板书边栏，避免四个未知量同时涌入一条长方程。

状态 C 应先画并联两端电压 \(U_3=I_3r_g\)，从而得到 \(R_0\)
支路电流，再写电池总电流。状态 D 同样从含表头支路电压
\(U_4=I_4(R_x+r_g)\)
出发。统一从``表头读数---端电压---旁路电流---总电流''推进，两个二次方程自然呈现同型结构。

求根时不直接说``取小根''，而应写三条筛选：接近给定估计值；所有电阻为正；回代四状态且不超量程。若题目删去估计值，单靠四读数可能存在交换根的离散二义性，需要额外状态或电压测量消除。

实际评分可以把电路图与代数各占一半。接线图必须能实现四状态且无危险短路；代数中必须先定义
\(S\)，说明 \(I_3,I_4\)
是支路电流，再列二次方程和根筛选。只给四个最终公式而没有说明每个电流读数的物理含义，不能视为完成``设计并测量''的要求。反之，电路构思正确但只写``联立可解''也缺少可复核性。

\subsection{九、变式训练}\label{ux4e5dux53d8ux5f0fux8badux7ec3-6}

若电流表量程不足，允许再增加一个已知保护电阻
\(R_p\)，设计其固定接法，使四种状态均不超量程，并写出修正后的方程。

最直接的保护方式是把 \(R_p\)
永久串在电池公共支路中，使四种状态的电池总电阻都增加
\(R_p\)，且不改变各并联支路内部的分流比。此时状态 A、B 变为

\[
\varepsilon=I_1(S+R_0+R_p),
\qquad
\varepsilon=I_2(S+R_p). \tag{29}
\]

先解的是 \(S+R_p\)，再减去已知 \(R_p\) 得 \(S\)。状态 C、D
的电池内阻项应统一替换为 \(r+R_p\) 对总电流的压降，但 \(R_p\)
不属于待测内阻，最终反演后需扣除。选择 \(R_p\)
时应按最小总电阻状态估计最大电流，保证表头支路电流小于 \(I_g\)；同时
\(R_p\) 不宜过大，否则各状态读数过于接近，反演误差放大。

另一种变式是把保护电阻只与表头串联，但这样会把有效表头内阻变为
\(r_g+R_p\)，所有含 \(r_g\)
的方程都要相应修正。训练重点是：保护元件放在公共支路还是测量支路，会以不同方式进入方程，不能只在最终答案中随意加一个
\(R_p\)。

\newpage

\section{第 8 题}\label{ux7b2c-8-ux9898}

质量可忽略的刚体棒两端固定质量为
\(m_1=1\,\mathrm{kg}\)、\(m_2=2\,\mathrm{kg}\) 的小球，铰点 \(O\)
到两端的距离分别为
\(r_1=0.4\,\mathrm m\)、\(r_2=0.6\,\mathrm m\)。棒绕竖直轴转动，初始时棒与竖直轴夹角
\(\theta_0=30^\circ\)，\(m_1\) 位于上半平面，\(m_2\)
位于下半平面，\(g=9.8\,\mathrm{m/s^2}\)。

（1）若两球以固定 \(\theta_0\) 作水平圆周运动，求角速度 \(\omega_0\)。

（2）若初始角速度 \(\omega_1=2\omega_0\)，求运动过程中 \(\theta\)
的最大值和最小值。

（3）另有一根总长为 \(l\) 的轻杆，两端质量均为
\(m\)，中点固连套筒并套在竖直轴上。套筒使夹角 \(\theta\)
保持不变，且以恒定角速度 \(\omega\)
转动。不得使用转动参考系，求杆和套筒系统对中心所受合外力矩的大小。

\begin{figure}[H]
\centering
\includegraphics[width=0.72\linewidth,height=\textheight,keepaspectratio,alt={图8-1　不对称双球杆与带套筒对称双球杆的空间转动}]{../figures/question-08.pdf}
\caption*{图8-1　不对称双球杆与带套筒对称双球杆的空间转动}
\end{figure}

\subsection{一、题型定位}\label{ux4e00ux9898ux578bux5b9aux4f4d-7}

锥角改变时，竖直轴角动量守恒而角速度改变。先用有效势找两个转折角，再在惯性系求对称双球杆的角动量矢量及其变化率。

\subsection{二、知识点}\label{ux4e8cux77e5ux8bc6ux70b9-7}

\begin{enumerate}
\def\labelenumi{\arabic{enumi}.}
\tightlist
\item
  定夹角匀速进动可由关于铰点、垂直于杆与轴平面的力矩平衡求角速度。
\item
  约束力对竖直轴无力矩，故 \(L_z\) 守恒。
\item
  理想铰接和重力场中机械能守恒。
\item
  空间转动中角动量一般不与角速度平行。
\item
  合外力矩满足
  \(\boldsymbol\tau=d\boldsymbol L/dt\)，必须在惯性系中求导。
\end{enumerate}

\subsection{三、物理图像}\label{ux4e09ux7269ux7406ux56feux50cf-7}

不对称双球杆的重心在 \(m_2\)
一侧，重力倾向于使重球向下；转动产生的几何效应可维持固定锥角。角速度改变后，锥角会在两个转折值之间摆动。对称双球杆虽重力矩相消，但角动量矢量随杆的方位转动，套筒仍需提供非零力矩。

\subsection{四、解题入口}\label{ux56dbux89e3ux9898ux5165ux53e3-7}

第一问直接写定锥角条件。第二问先用 \(L_z\) 守恒消去任意时刻的
\(\omega\)，再代入机械能守恒求转折角。第三问在惯性系中写出两球的总角动量矢量并求时间导数。

\subsection{五、解析}\label{ux4e94ux89e3ux6790-7}

\begin{figure}[H]
\centering
\includegraphics[width=0.88\linewidth,height=\textheight,keepaspectratio,alt={图8-2　锥角、角动量分量与套筒约束力矩}]{../figures/solution-08.pdf}
\caption*{图8-2　锥角、角动量分量与套筒约束力矩}
\end{figure}

图8-2中的 \(\bm L\) 不沿杆轴，\(L_z\) 保持不变；轴力矩垂直于子午面，指向
\(-\bm e_\phi\)。它改变角动量方向，且与竖直角速度垂直。

\needspace{3\baselineskip}
\textbf{（一）固定锥角运动的力矩平衡}

定义

\[
K=m_1r_1^2+m_2r_2^2. \tag{1}
\]

两球到竖直轴的水平距离分别为
\(r_1\sin\theta_0\)、\(r_2\sin\theta_0\)。定锥角转动时，惯性系中每个球具有指向轴线的向心加速度。

把各力对 \(O\) 的力矩投影到垂直于``杆---竖直轴''平面的方向，可得

\[
\left[
K\omega_0^2\cos\theta_0
-(m_2r_2-m_1r_1)g
\right]\sin\theta_0=0. \tag{2}
\]

\(\theta_0\ne0\)，所以

\[
K\omega_0^2\cos\theta_0
=(m_2r_2-m_1r_1)g. \tag{3}
\]

因此

\[
\boxed{
\omega_0=
\sqrt{
\frac{(m_2r_2-m_1r_1)g}
{(m_1r_1^2+m_2r_2^2)\cos\theta_0}
}
}. \tag{4}
\]

代入数据：

\[
m_2r_2-m_1r_1
=2\times0.6-1\times0.4
=0.8\,\mathrm{kg\,m}, \tag{5}
\]

\[
K=1\times0.4^2+2\times0.6^2
=0.88\,\mathrm{kg\,m^2}. \tag{6}
\]

故

\[
\begin{aligned}
\omega_0
&=\sqrt{
\frac{0.8\times9.8}
{0.88\cos30^\circ}
}\\
&\simeq3.21\,\mathrm{rad/s}.
\end{aligned} \tag{7}
\]

\needspace{3\baselineskip}
\textbf{（二）竖直轴角动量守恒}

任意时刻，关于竖直轴的转动惯量为

\[
I_z=K\sin^2\theta. \tag{8}
\]

没有外力对竖直轴的力矩，因此

\[
I_z\omega=\text{常量}. \tag{9}
\]

初态 \(\omega_1=2\omega_0\)，所以

\[
K\sin^2\theta\,\omega
=K\sin^2\theta_0\,\omega_1. \tag{10}
\]

约去 \(K\)：

\[
\boxed{
\omega
=\omega_1
\frac{\sin^2\theta_0}{\sin^2\theta}
}. \tag{11}
\]

\needspace{3\baselineskip}
\textbf{（三）机械能守恒求转折角}

转折位置有 \(\dot\theta=0\)。此时动能只有绕竖直轴的转动：

\[
T=\frac12K\omega^2\sin^2\theta. \tag{12}
\]

取 \(O\) 为重力势能零高度基准。\(m_1\) 高度为 \(r_1\cos\theta\)，\(m_2\)
高度为 \(-r_2\cos\theta\)，故

\[
U=(m_1r_1-m_2r_2)g\cos\theta. \tag{13}
\]

初态与任一转折态能量相等：

\[
\frac12K\omega^2\sin^2\theta
+(m_1r_1-m_2r_2)g\cos\theta
=
\frac12K\omega_1^2\sin^2\theta_0
+(m_1r_1-m_2r_2)g\cos\theta_0. \tag{14}
\]

代入式（11）：

\[
\frac12K\omega_1^2
\frac{\sin^4\theta_0}{\sin^2\theta}
-(m_2r_2-m_1r_1)g\cos\theta
=
\frac12K\omega_1^2\sin^2\theta_0
-(m_2r_2-m_1r_1)g\cos\theta_0. \tag{15}
\]

再利用第一问关系

\[
(m_2r_2-m_1r_1)g
=K\omega_0^2\cos\theta_0, \tag{16}
\]

以及 \(\omega_1=2\omega_0\)，约去 \(K\omega_0^2\) 后得到

\[
\frac{2\sin^4\theta_0}{\sin^2\theta}
-\cos\theta_0\cos\theta
=2\sin^2\theta_0-\cos^2\theta_0. \tag{17}
\]

代入 \(\theta_0=30^\circ\)，令 \(c=\cos\theta\)：

\[
\frac{1}{8(1-c^2)}
-\frac{\sqrt3}{2}c
=-\frac14. \tag{18}
\]

两端乘以 \(8(1-c^2)\)：

\[
1-4\sqrt3c(1-c^2)=-2(1-c^2). \tag{19}
\]

展开并移到左边：

\[
4\sqrt3c^3-2c^2-4\sqrt3c+3=0. \tag{20}
\]

其中初始位置 \(c_0=\cos30^\circ=\sqrt3/2\) 必为一根，因此因式分解为

\[
4\sqrt3
\left(c-\frac{\sqrt3}{2}\right)
\left(c^2+\frac{\sqrt3}{3}c-\frac12\right)=0. \tag{21}
\]

二次因子的两个根为

\[
c=\frac{-\sqrt3\pm\sqrt{21}}{6}. \tag{22}
\]

带负号的根小于 \(-1\)，不可能是余弦值；带正号的根

\[
c=\frac{\sqrt{21}-\sqrt3}{6}
\simeq0.47509
\]

对应

\[
\theta=\arccos c\simeq61.6^\circ.
\]

所以在 \(0<\theta<\pi/2\) 内的两个转折位置为

\[
\theta=30.0^\circ,
\qquad
\theta\simeq61.6^\circ. \tag{23}
\]

由于初态就是一个转折位置，且随后锥角增大，所以

\[
\boxed{
\theta_{\min}=30.0^\circ,\qquad
\theta_{\max}\simeq61.6^\circ
}. \tag{24}
\]

\needspace{3\baselineskip}
\textbf{（四）对称双球杆的角动量矢量}

设 \(\hat{\boldsymbol z}\) 为竖直单位矢量，\(\boldsymbol n\)
为沿杆从中心指向上端小球的单位矢量。两个小球的位置为

\[
\boldsymbol r_\pm
=\pm\frac l2\boldsymbol n. \tag{25}
\]

杆以角速度

\[
\boldsymbol\omega=\omega\hat{\boldsymbol z} \tag{26}
\]

转动。速度为

\[
\boldsymbol v_\pm
=\boldsymbol\omega\times\boldsymbol r_\pm. \tag{27}
\]

单个质点角动量

\[
\boldsymbol L_\pm
=\boldsymbol r_\pm\times m\boldsymbol v_\pm. \tag{28}
\]

使用三重积恒等式

\[
\boldsymbol r\times(\boldsymbol\omega\times\boldsymbol r)
=r^2\boldsymbol\omega
-(\boldsymbol r\cdot\boldsymbol\omega)\boldsymbol r, \tag{29}
\]

两球相加：

\[
\boxed{
\boldsymbol L
=\frac{ml^2}{2}\omega
\left(
\hat{\boldsymbol z}-\cos\theta\,\boldsymbol n
\right)
}. \tag{30}
\]

这里 \(\boldsymbol L\) 并不平行于 \(\hat{\boldsymbol z}\)。

\needspace{3\baselineskip}
\textbf{（五）在惯性系中求导}

\(\omega\)、\(\theta\) 均为常量，但 \(\boldsymbol n\) 绕竖直轴转动：

\[
\frac{d\boldsymbol n}{dt}
=\boldsymbol\omega\times\boldsymbol n
=\omega\hat{\boldsymbol z}\times\boldsymbol n. \tag{31}
\]

对式（30）求导：

\[
\begin{aligned}
\boldsymbol\tau_{\rm ext}
=\frac{d\boldsymbol L}{dt}
&=-\frac{ml^2}{2}\omega\cos\theta
\frac{d\boldsymbol n}{dt}\\
&=-\frac{ml^2}{2}\omega^2\cos\theta
(\hat{\boldsymbol z}\times\boldsymbol n).
\end{aligned} \tag{32}
\]

因为

\[
|\hat{\boldsymbol z}\times\boldsymbol n|=\sin\theta, \tag{33}
\]

所以合外力矩大小

\[
\boxed{
\tau_{\rm ext}
=\frac12m\omega^2l^2|\sin\theta\cos\theta|
}. \tag{34}
\]

两端重力对中心的力矩互相抵消，因此这一力矩由套筒和轴的约束提供；杆与套筒系统对中心所受的合外力矩即为式（34）。

\needspace{3\baselineskip}
\textbf{补充推导与适用条件}

前两问可统一写成锥角的一维有效势问题。关于竖直轴的守恒角动量记为

\[
L_z=K\sin^2\theta\,\omega. \tag{35}
\]

总动能包含锥角变化和绕轴转动：

\[
T=\frac12K\dot\theta^2
+\frac12K\sin^2\theta\,\omega^2. \tag{36}
\]

消去 \(\omega=L_z/(K\sin^2\theta)\)，机械能写成

\[
E=\frac12K\dot\theta^2+V_{\rm eff}(\theta), \tag{37}
\]

其中

\[
V_{\rm eff}(\theta)
=\frac{L_z^2}{2K\sin^2\theta}
-(m_2r_2-m_1r_1)g\cos\theta. \tag{38}
\]

第一项是固定 \(L_z\) 后的``离心势垒''，当 \(\theta\to0\)
时发散，阻止系统在非零 \(L_z\)
下接近竖直轴；第二项是重力势能。定锥角运动就是
\(V_{\rm eff}'(\theta_0)=0\)，求导恰好恢复式（3）。因此第一问的力矩平衡与有效势极值是同一条件。

在第二问初态 \(\omega_1=2\omega_0\)。用

\[
L_z=2K\omega_0\sin^2\theta_0
\]

代入式（38），在 \(\theta_0\) 处求导：

\[
V_{\rm eff}'(\theta_0)
=-3K\omega_0^2\sin\theta_0\cos\theta_0<0. \tag{39}
\]

由 \(K\ddot\theta=-V_{\rm eff}'\)，初始
\(\ddot\theta>0\)，所以锥角先增大。这为式（24）中``\(30^\circ\)
是最小值、\(61.6^\circ\)
是最大值''提供了方向依据，而不是凭数值大小猜测。

式（18）的两个根必须回代原能量方程。\(\theta=30^\circ\)
是已知初始转折点；另一根约 \(61.6^\circ\)。还要检查
\(0<\theta<\pi/2\)、\(\sin\theta\ne0\)，以及有效势在两根之间不高于总能量。若整理过程中乘了
\(1-c^2\)，代数上可能引入 \(c=\pm1\) 的伪根，必须排除。

第三问的角动量可以分成竖直与沿杆方向的组合。式（30）中
\(\hat{\boldsymbol z}\) 不变，而 \(\boldsymbol n\)
绕轴转动，故只有第二项需要求导。角动量的水平分量大小为

\[
L_\perp=\frac{ml^2}{2}\omega\sin\theta\cos\theta. \tag{40}
\]

这个水平分量以角速度 \(\omega\) 转动，所以所需力矩大小也可直接写成

\[
\tau=\omega L_\perp
=\frac12m\omega^2l^2\sin\theta\cos\theta, \tag{41}
\]

与式（34）一致。它提供比完整矢量求导更直观的二次检查。

极限同样重要：\(\theta=0\)
时两球都在轴上，没有圆周运动，力矩为零；\(\theta=\pi/2\)
时杆水平，角动量沿竖直轴，也不需改变方向，力矩仍为零；中间在
\(45^\circ\) 达到最大。这些结论与 \(\sin\theta\cos\theta\) 完全一致。

\needspace{3\baselineskip}
\textbf{初态条件、允许区间与约束功}

第二问使用初始
\(\dot\theta=0\)：题目的``初始绕轴转动''按参考解理解为没有额外的锥角方向速度。若另给非零
\(\dot\theta\)，总能量要加
\(K\dot\theta^2/2\)，两个转折角将改变。图8-2画出的有效势对应这一初态，并在
\(30^\circ\) 与 \(61.6^\circ\) 之间标出 \(E\ge V_{\rm eff}\)
的允许区间。

第三问中 \(L_z=(ml^2/2)\omega\sin^2\theta\) 恒定，水平角动量却为
\(\boldsymbol L_\perp=-(ml^2/2)\omega\sin\theta\cos\theta\,\hat{\boldsymbol e}_\rho\)，随方位转动。因此

\[
\boldsymbol\tau=-\frac{ml^2}{2}\omega^2\sin\theta\cos\theta\,
\hat{\boldsymbol e}_\varphi.
\]

轴与套筒提供水平约束力矩，且
\(\boldsymbol\tau\cdot\omega\hat{\boldsymbol z}=0\)，不改变转动动能。这与套筒可以绕轴自由转动一致：自由转动只意味着不传递绕竖直轴的力矩，并不意味着各方向约束力矩均为零。若允许
\(\theta>\pi/2\)，所求大小应写绝对值；在题图的锐角范围内可去掉绝对值。

\subsection{六、关键洞察}\label{ux516dux5173ux952eux6d1eux5bdf-7}

解析图中的 \(\bm L\) 不沿杆轴，\(L_z\)
保持不变；轴力矩垂直于子午面，指向
\(-\bm e_\phi\)。它改变角动量方向，且与竖直角速度垂直。

\subsection{七、易错点}\label{ux4e03ux6613ux9519ux70b9-7}

\begin{enumerate}
\def\labelenumi{\arabic{enumi}.}
\tightlist
\item
  第一问把 \(m_2\) 的重力矩方向写反。
\item
  第二问错误地认为角速度守恒；守恒的是 \(L_z\)。
\item
  势能中漏掉 \(m_2\) 位于铰点下方的负高度。
\item
  只由机械能守恒求 \(\theta\)，没有同时使用角动量守恒。
\item
  第三问把 \(\boldsymbol L\) 误写成
  \(I_z\omega\hat{\boldsymbol z}\)，漏掉水平分量。
\item
  第二问只用 \(L_z\) 守恒却漏掉
  \(\tfrac12K\dot\theta^2\)；虽然转折点该项为零，判断全过程仍需要它。
\item
  整理转折方程时乘以 \(\sin^2\theta\) 后保留 \(\theta=0\) 的伪根。
\item
  把``合外力矩对中心''误读为关于竖直轴的力矩；第三问所求矢量力矩水平且随杆转动。
\item
  认为两端重力矩相消就表示总外力矩为零，忽略套筒和轴的约束力矩。
\end{enumerate}

\subsection{八、课堂讲解建议}\label{ux516bux8bfeux5802ux8bb2ux89e3ux5efaux8bae-7}

先在图上标出两球的水平圆半径和高度，分别建立 \(I_z(\theta)\) 与
\(U(\theta)\)。第三问可用一支倾斜旋转的雨伞骨架作类比：姿态角不变，但杆方向持续改变，因此角动量方向也持续改变。

图 8-1 必须按三维关系阅读：杆与竖直轴夹角是
\(\theta\)。两球的水平圆半径分别为

\[
r_1\sin\theta,\qquad r_2\sin\theta,
\]

高度分别为

\[
+r_1\cos\theta,\qquad -r_2\cos\theta.
\]

课堂先让学生从图上写出这些投影，再进入力矩；若三维图只当装饰，后续
\(I_z\) 与势能符号很容易错。

第一问可以先对每个小球写向心加速度，再把合力矩投影到垂直于``杆---轴''平面的方向。推得式（3）后用单位检查：分子是
\(\mathrm{kg\,m^2/s^2}\)，分母是 \(\mathrm{kg\,m^2}\)，所以
\(\omega_0^2\) 单位正确。

第二问板书分两条守恒律。先由无竖直轴外力矩得
\(L_z\)，再由理想铰链不做功得机械能。把 \(\omega(\theta)\)
代入能量后才求转折角。求出两个根后用式（39）判断初始运动方向，并在有效势图上标出允许区间。

第三问严格遵守题目``不使用转动参考系''：只写惯性系位置、速度、角动量和时间导数。可先让学生猜``角速度不变是否等于力矩为零''，再用式（40）说明水平角动量不断转向。验收时必须看到矢量三重积或等价的水平分量转动论证。

教师可把第三问与匀速圆周运动类比：速率不变但速度方向改变，需要向心力；同理，角动量相关分量大小不变但方向改变，需要``使角动量转向的力矩''。水平角动量端点在水平面画圆，其变化率大小是
\(\omega L_\perp\)。这个几何图像比直接背
\(\boldsymbol\tau=d\boldsymbol L/dt\)
更容易迁移到陀螺进动和非主轴转动问题。

\subsection{九、变式训练}\label{ux4e5dux53d8ux5f0fux8badux7ec3-7}

若对称双球杆的夹角 \(\theta\) 允许缓慢变化，写出总动能

\[
T=\frac{ml^2}{4}\left(\dot\theta^2+\omega^2\sin^2\theta\right),
\]

并讨论竖直轴角动量守恒如何决定 \(\omega(\theta)\)。

两球对中心的转动惯量系数为 \(K_s=ml^2/2\)。题给总动能可写成

\[
T=\frac12K_s\dot\theta^2
+\frac12K_s\omega^2\sin^2\theta. \tag{42}
\]

由于关于竖直轴无外力矩，

\[
L_z=K_s\omega\sin^2\theta=\text{常量},
\]

所以

\[
\omega(\theta)=\frac{L_z}{K_s\sin^2\theta}. \tag{43}
\]

\Needspace{7\baselineskip}

夹角减小时角速度迅速增大。代回动能得到有效势项
\(L_z^2/(2K_s\sin^2\theta)\)。若另有使 \(\theta\)
改变的势能或弹性约束，只需与该有效势相加即可判断转折点和小振动。训练中要强调：\(\omega\)
不是守恒量，它由 \(L_z\) 守恒和几何转动惯量共同决定。
\end{document}
